% Modernised reading — fonds Alexandre Grothendieck, shelfmark 64, pages 1-156
% (the whole folder, eight batches of transcription).
%
% This is an INTERPRETATION, not a transcription. It re-reads the pages in
% today's notation and vocabulary, states the hypotheses the manuscript leaves
% implicit, and drops the critical apparatus so that the mathematics reads
% straight through. Everything the brackets and underlines carried — a doubtful
% reading, a notation he used differently, a missing passage — moves into a
% footnote. It is held to being mathematically correct as it stands; where the
% pages are loose or mistaken, the true statement is what appears here and the
% footnote says what the page has. In any dispute of substance the
% transcription governs: whoever cites Grothendieck must cite the
% batch-NN.fr.tex files.
%
% Scope: all eight batches (pages 1-156) were read whole before any of this
% was written; ONE file for the shelfmark. Pages without mathematics (40, 123,
% 126: typed administrative documents; 121, 129, 141, 152: blank) are not read.
%
% What the folder is. Five parts in his hand, each with its own cover or
% heading and its own pagination: I « Fourbis combinatoires » (pp. 1-42),
% II « Structures octaédrales dans P1 » (pp. 43-68), III « Épinglages de
% Legendre » (pp. 69-99), IV « Descriptions explicites par fonctions P de
% Weierstrass » (pp. 100-125), « Approche transcendante » (pp. 128-151, in two
% drafts), and an arithmetic appendix (pp. 153-156). A heading « Relations
% entre M11 et M04 » (p. 127) has no text under it.
% The spine: the 4-torsion of an elliptic curve in char != 2, M = 4E, a free
% rank-2 Z/4-module, is described combinatorially (Part I: M <-> (I, Q, k),
% and the octahedron Delta_M = P(M)(Z/4)), then geometrically (Part II: the
% 6 classes of points of order 4 mod +-1 form an octahedral configuration in
% E/+-1 = P^1), then used to rigidify E (Part III: Legendre pinnings, fine
% moduli U_{0,3} x mu_4^*), then made explicit (Part IV) and over C (V).
%
% Notation fixed for the whole document. X^* for a finite set or group X is
% his « ensemble des éléments non nuls » (V^* = V \ {0}), NEVER the dual; the
% dual is X^vee. ω(X) (underlined omega) the 2-element set of orientations.
% His « cube » of Part I §2 (four vertices, four edges, group of order 8) is a
% SQUARE, the 2-dimensional cube; it is called « carré » here, as he himself
% calls it in Part III. « Cube » is kept for the 3-dimensional cube of §4.
% D_4 = dihedral group of order 8 (his \mathbb{D}_4); the Klein four-group is
% written (Z/2Z)^2 (he writes D_2 or D_4 for it in Part II). Contracted
% products of torsors are written with ∧. Part IV: y = -(1/2) Θ_δ ℘ replaces
% his ℘' (see below). Elements of μ_4^* are written ι to keep i for J.
%
% Substantive departures from the pages, each footnoted where it occurs:
% — p. 3: « pleinement fidèle » holds for the groupoids of isomorphisms only.
% — p. 6: x -> Cent_G(x) lands on the 2-Sylows, not on the cyclic groups of
%   order 4 (the second diagram of the page has it right).
% — p. 9: D^+_ab is D_ab; p. 10: the enumeration « B, S, ω » is D, C, ω_Q;
%   p. 20 d): V, not V^*; p. 21: 19 (not 17) subgroups of 2-power order.
% — p. 13: ϖcd = 1 holds with products read as right actions; with the usual
%   composition the relation is dcϖ = 1 (checked, ours).
% — p. 16: the functor (groups ≅ D_4) -> (squares) IS an equivalence (the page
%   is illegible on whether it says so); verified, ours.
% — p. 31: GL(M) = (GL(M_0)·M_0)·F_2^ω is semidirect, not direct, in its
%   second factor.
% — p. 41: V' = Aut^-(V) ∪ {0} is not a commutative subalgebra (σ_iσ_j ≠
%   σ_jσ_i); W = Aut^+(V) ∪ {0} is, and is the field F_4.
% — pp. 49-52: the theorem on octahedral configurations needs the residue
%   characteristics ≠ 5 as well as ≠ 2 wherever it says Aut(X, Δ) is the
%   octahedral group: in char 5, Δ = P^1(F_5) and Aut(X, Δ) = PGL_2(F_5).
%   The bijections a)-b)-c) themselves survive with Δ carrying its antipodism.
% — pp. 51, 52, 59: « icosaédrale » read octaédrale.
% — p. 63: Δ^t_∞ = {t ± sqrt(t(t-1))}, not sqrt(t(1-t)).
% — p. 102: the right-hand map of (9) is multiplication by -nδ, not δ; hence
%   Ker Θ_δ = k only in char 0 (in char p it contains the p-th powers).
% — pp. 104-107: ε_3(Θ_δ ℘) = -2δ^3, so (Θ_δ ℘)^2 = 4(℘^3 + ...); the
%   page's monic relation (28) and everything after hold for y = -Θ_δ℘/2.
% — p. 117: the bracket in (71) and the margin carry sign slips; the clean
%   form (Z_i - ζ_i)^2 = (ζ_i - ζ_j)(ζ_i - ζ_k) is given (ours).
% — p. 138: the family E_z = C/(Z + Z(-z̄)) is anti-holomorphic in z; with
%   (24) it is C/(Z + Zz). p. 139: H = π_ab has rank 2.
% — p. 154: « B normal ssi a ∉ m^2 » is true for A a DVR (the case used:
%   normality in codimension 1), false in dimension >= 2.
%
% Announced and not established: the inclusion diagram of the 30 subgroups
% of S_4 (p. 21); the cohomological lifting of an octahedron to a Z/4-module
% (pp. 36-37); an intrinsic definition of f^t_2 (p. 62); the « détermination
% de M_11(4) » (pp. 65-68, stops mid-sentence); the equivalence of the Kummer
% digression (p. 73, « je dis »); the canonical Legendre pinning defined by a
% level-4 Jacobi pinning (pp. 92, 117-120: « Dieu sait comment --- il faudra y
% revenir »); condition (73) in coordinates (p. 99); « Relations entre M11 et
% M04 » (p. 127, a title); the universal cover of PD^+ (pp. 140, 151); the
% sufficiency claimed on p. 155 and the example after it.
%
% Pass: Opus 5.5 (claude-opus-5-5), 2026-10-03 — first pass, unchecked against the pages by a human.
\documentclass[11pt,a4paper]{article}
\input{../preamble/grothendieck}

\folder{64}
\pages{1}{156}
\dating{[à partir de 1980- 1982]}
\watermark{Édition de démonstration\\interprétation personnelle de l'œuvre}
\foldertitle{Modules courbes elliptiques en caractéristiques ≠2 : notes manuscrites (s.d.) — lecture modernisée du dossier entier}

\begin{document}

\section*{Résumé}

\begin{resume}
Une courbe elliptique est, vue de loin, un tore : une surface en forme de
bouée, sur laquelle on sait additionner les points. Certains points y jouent
un rôle à part, ceux qui reviennent à l'origine quand on les ajoute deux fois à
eux-mêmes (il y en a quatre), ou quatre fois (il y en a seize). Ce dossier,
écrit au début des années 1980, se demande ce que ces seize points
\emph{savent} de la courbe, et comment s'en servir pour la décrire sans rien
choisir.

Le problème est un problème de catalogue. On voudrait un espace dont chaque
point représente une courbe elliptique, et une seule. Mais toute courbe
elliptique possède une symétrie, le retournement $x \mapsto -x$, et ces
symétries empêchent le catalogue d'exister honnêtement : c'est comme vouloir
ranger des papillons sans savoir de quel côté les épingler. On s'en sort en
ajoutant aux courbes une petite structure qui tue les symétries ---
c'est ce qu'il appelle un \emph{épinglage}, et le mot est juste : on pique la
courbe pour qu'elle ne puisse plus se retourner.

La première idée du dossier est que les seize points cachent des polyèdres
réguliers. Les quatre points tués par deux forment un ensemble à quatre
éléments, et un tel ensemble, sans rien d'autre, porte déjà une géométrie
cachée : ses permutations sont exactement les rotations d'un cube. Les douze
points d'ordre exactement quatre, pris au signe près, forment six points sur
une droite projective --- et ces six points sont les sommets d'un octaèdre,
aussi rigidement que $\{0, \infty, \pm 1, \pm i\}$ sur la sphère de Riemann.
La première partie fait l'inventaire, purement combinatoire, de toutes ces
correspondances ; la deuxième montre qu'elles se réalisent géométriquement.

La seconde idée est l'épinglage de Legendre. Pour supprimer le retournement, il
faut choisir un vecteur tangent à l'origine --- mais au signe près seulement il
est déterminé par la courbe, et le signe est précisément ce que le retournement
change. Grothendieck montre qu'on lève l'ambiguïté en payant un prix
arithmétique exactement mesuré : le choix d'une racine carrée de $-1$. Le
catalogue des courbes munies de cet épinglage et d'une numérotation de leurs
points d'ordre deux existe alors vraiment : c'est la droite privée de trois
points, doublée par les deux choix de $i$. Les dernières parties rendent tout
explicite, par les fonctions de Weierstrass, puis sur les nombres complexes par
le demi-plan de Poincaré.

Une manière se lit à chaque page : rien n'est jamais « égal », tout est
« canoniquement isomorphe », et chaque fois qu'un isomorphisme dépend d'un
choix il le dit, compte les choix, et cherche lequel imposer. « Il faut
expliciter de bout en bout la convention », écrit-il ; puis, quand une
convention lui donne enfin la formule de tout le monde : « victoire, c'est la
formule habituelle » ; et ailleurs, devant un calcul qu'il ne voit pas encore :
« Dieu sait comment --- il faudra y revenir ». Les pages sont contemporaines de
la \emph{Longue Marche à travers la théorie de Galois} (1981), où les espaces
de modules de courbes et leurs groupes fondamentaux sont au centre ; un titre
du dossier, resté sans texte, annonce les « relations entre $M_{1,1}$ et
$M_{0,4}$ ».

\keywords{Klein four-group, symmetric group S4, dihedral group D4, octahedral group, contracted product of torsors, GL2(Z/4Z), level-4 structure, 4-torsion points, octahedral configuration, Klein quotient of P1, moduli of four points on P1, Legendre family, lambda-line, rigidification, tamely ramified Kummer covering, tangential base point, fine moduli space, dicyclic group, Weierstrass function, chord-tangent addition law, division points, upper half-plane, SL2(Z), mapping class group of the torus, Teichmüller group, normality of quadratic extensions}
\end{resume}

\section*{Le fil du dossier, et les conventions}

Le dossier se compose de parties que Grothendieck a lui-même séparées par des
feuillets de titre, chacune avec sa pagination propre.

\begin{enumerate}
\item[I.] \emph{Fourbis combinatoires} (pages 1 à 42). Les ensembles à quatre
éléments et le groupe $\mathfrak{S}_4$ (pages 2 à 6) ; le carré combinatoire et
son groupe diédral d'ordre $8$ (pages 7 à 19) ; les sous-groupes de
$\mathfrak{S}_4$ (pages 19 à 21) ; le cube et l'octaèdre attachés à un ensemble
à quatre éléments (pages 22 à 24) ; les $\mathbb{Z}/4\mathbb{Z}$-modules libres
de rang $2$ (pages 25 à 31) ; l'octaèdre d'un tel module et une récapitulation
(pages 32 à 39) ; une proposition sur $\mathrm{End}(V)$, $V$ plan sur
$\mathbb{F}_2$ (pages 41 et 42).
\item[II.] \emph{Structures octaédrales dans $\mathbb{P}^1$} (pages 43 à 68) :
les actions du groupe de Klein sur un fibré en droites projectives et la
configuration octaédrale ; les courbes de type $(0,4)$ ; la relation avec les
courbes elliptiques, et l'application $f_2^t$ induite par la multiplication
par $2$ ; un début de « détermination de $M_{1,1}(4)$ ».
\item[III.] \emph{Épinglages de Legendre} (pages 69 à 99) : une digression sur
les revêtements de Kummer épinglés par les espaces tangents ; l'épinglage de
Legendre sur un corps, puis sur une base ; le schéma de modules fin
$U_{0,3} \times \mu_4^{*}$ et son groupe d'automorphismes ; une explicitation
en coordonnées.
\item[IV.] \emph{Descriptions explicites par fonctions $\wp$ de Weierstrass}
(pages 100 à 125) : l'algèbre graduée des fonctions à pôle à l'origine,
l'équation de Weierstrass, la loi d'addition, les points d'ordre $4$, et un
début de traduction en épinglage de Legendre.
\item[V.] \emph{Approche transcendante} (pages 127 à 151), en deux
\emph{moutures} qui ne se recouvrent pas tout à fait : une première (pages
130 à 140), qui passe par les structures complexes sur un plan réel, et une
seconde (pages 142 à 151), plus directe. On les garde distinctes.
\item[VI.] Un appendice arithmétique (pages 153 à 156) : normalité des
extensions quadratiques, revêtements quadratiques de $\mathbb{P}^1_{\mathbb{Z}}$.
\end{enumerate}

\emph{Le fil.} Tout tourne autour d'un seul objet, le module $M = {}_4E$ des
points de $4$-torsion d'une courbe elliptique en caractéristique résiduelle
$\neq 2$, qui est localement un $\mathbb{Z}/4\mathbb{Z}$-module libre de rang
$2$. La partie I le décrit combinatoirement : un tel module équivaut à un
ensemble à quatre éléments, un carré et une identification entre deux
ensembles à deux éléments (page 28), et il porte un octaèdre,
$\Delta_M = \mathbb{P}(M)(\mathbb{Z}/4\mathbb{Z})$ (page 32). La partie II
montre que cet octaèdre est réalisé sur la courbe : les points d'ordre $4$ au
signe près forment, dans $E/\{\pm 1\} \simeq \mathbb{P}^1$, une configuration
octaédrale (page 60). La partie III s'en sert pour épingler $E$ ; la partie IV
et l'approche transcendante rendent les choses explicites.

\emph{Deux homonymies, fixées ici.} (1) Le « cube combinatoire » de la partie I,
§2, a quatre sommets, quatre arêtes et un groupe d'automorphismes d'ordre $8$ :
c'est le cube de dimension $2$, un \emph{carré}. On l'appelle ainsi ; c'est
d'ailleurs le mot de Grothendieck dans la partie III (« un carré
combinatoire », page 75).\footnote{Les figures marginales des pages 12 et 13
sont des carrés, et une note marginale de la page 16 dit « carré ».} Le mot
« cube » est réservé au cube de dimension $3$ que construit le §4. (2) La lettre
$\mathbb{D}_4$ désigne le groupe diédral d'ordre $8$, groupe du carré. Le
groupe de Klein, que la partie II note tantôt $\mathbb{D}_2$, tantôt
$\mathbb{D}_4$, est noté ici $(\mathbb{Z}/2\mathbb{Z})^2$.

\emph{Conventions.} Pour un ensemble ou un groupe abélien $X$, $X^{*}$ désigne,
comme sur les pages, l'ensemble des éléments \emph{non nuls} (ou $\neq 1$) ; le
dual est noté $X^{\vee}$. Pour un ensemble fini $X$ de cardinal $\geq 2$,
$\underline{\omega}(X)$ est l'ensemble à deux éléments de ses orientations,
torseur sous $\{\pm 1\} = \mathfrak{S}_{X,\mathrm{ab}}$. Un ensemble à deux
éléments est un torseur sous $\mathbb{F}_2$ ; le produit contracté de deux tels
torseurs est noté $P \wedge P'$, et $\mathbb{1}$ est le torseur trivial.
Pour un groupe abélien $A$ et un ensemble fini $J$,
$V_J(A) = \operatorname{Ker}(A^J \xrightarrow{\text{somme}} A)$.
$\mathfrak{S}_I^{+}$ est le groupe alterné. Toute la géométrie se fait sur des
bases où $2$ est inversible.

\pagerange{1}{6}
\section*{I. Fourbis combinatoires --- 1. Ensembles à quatre éléments (pages 1 à 6)}

Soit $I$ un ensemble à quatre éléments. Il porte canoniquement une structure
de plan affine sur $\mathbb{F}_2$ : le sous-groupe $V \subset \mathfrak{S}_I$
formé de l'identité et des trois doubles transpositions est isomorphe --- non
canoniquement --- à $\mathbb{F}_2^2$ et opère sur $I$ de façon simplement
transitive. C'est le seul sous-groupe distingué abélien non trivial de
$\mathfrak{S}_I$, il est son propre centralisateur, et l'action par
conjugaison donne
\[
\mathfrak{S}_I / V \xrightarrow{\;\sim\;} \operatorname{Aut}_{\mathbb{F}_2}(V)
\xrightarrow{\;\sim\;} \mathfrak{S}_J, \qquad J = V^{*},
\]
de sorte que $1 \to V \to \mathfrak{S}_I \to \mathfrak{S}_J \to 1$ s'identifie
à la suite $1 \to V \to \operatorname{Aut}_{\mathrm{aff}}(I) \to
\operatorname{GL}(V) \to 1$ du groupe affine. Les trois éléments de $J$
correspondent aux trois partitions de $I$ en deux paires ; l'application
$A \mapsto \{A, I \setminus A\}$, des six parties à deux éléments vers ces
trois partitions, associe à une droite affine de $I$ sa direction.

Le foncteur d'oubli des plans affines sur $\mathbb{F}_2$ vers les ensembles à
quatre éléments est une équivalence. Plus intéressant : le foncteur
$I \mapsto \mathfrak{S}_I$, du groupoïde des ensembles à quatre éléments vers le
groupoïde des groupes isomorphes à $\mathfrak{S}_4$ et de leurs isomorphismes,
est une équivalence.\footnote{La page dit « pleinement fidèle », de la
catégorie $(\mathrm{Ens}_4)$ vers $(\mathrm{Gr})$. Entendu dans la catégorie de
tous les groupes et de tous les homomorphismes, c'est faux : il existe des
homomorphismes $\mathfrak{S}_I \to \mathfrak{S}_{I'}$ qui ne sont pas induits
par une bijection (le trivial, ou la signature suivie d'une inclusion). L'énoncé
vrai porte sur les groupoïdes d'isomorphismes, et c'est ce qu'il en tire :
$\mathfrak{S}_4$ est sans centre et tous ses automorphismes sont intérieurs.} Un
quasi-inverse associe à $G$ l'ensemble $I_G$ de ses quatre $3$-sous-groupes de
Sylow. En effet, ceux-ci correspondent aux scindages de l'extension
$1 \to V \to G^{+} \to \mathfrak{S}_J^{+} \to 1$ du groupe cyclique d'ordre $3$
par $V$ ; $V$ opère transitivement sur ces scindages par conjugaison, et
librement puisque $V^{\mathfrak{S}_J^{+}} = 0$ ; donc $I_G$ est un torseur sous
$V$. Pour $G = \mathfrak{S}_I$, la bijection $I \to I_G$ envoie $i$ sur le
$3$-Sylow qui fixe $i$ ; dans l'autre sens, l'action de $G$ sur $I_G$ par
conjugaison donne un isomorphisme $G \simeq \mathfrak{S}_{I_G}$.

\emph{Les éléments de $\mathfrak{S}_I$, et le cube.} Grothendieck range les
$24$ éléments en cinq classes et note en marge, pour chacune, son nom dans le
groupe des rotations du cube :
\begin{enumerate}
\item[a)] l'identité ;
\item[b)] les trois éléments de $J = V^{*}$, indexés par les partitions de type
$(2,2)$ --- les demi-tours autour des axes des faces ;
\item[c)] les huit éléments d'ordre $3$, qui s'identifient à
$I \times \underline{\omega}$, $\underline{\omega} = \underline{\omega}(J)$, par
leur point fixe et leur image dans $\mathfrak{S}_J^{+}$ ; ils forment deux
classes à droite et à gauche sous $V$ --- les rotations autour des grandes
diagonales ;
\item[d)] les six transpositions, indexées par les parties à deux éléments ;
elles se groupent en trois paires $\{\tau_A, \tau_{I \setminus A}\}$, de produit
dans $V^{*}$ --- les demi-tours autour des axes joignant les milieux d'arêtes
opposées ;
\item[e)] les six éléments d'ordre $4$, qui correspondent aux six ordres
circulaires sur $I$ --- les quarts de tour autour des axes des faces.
\end{enumerate}
Pour $u$ d'ordre $4$, $u^2 \in V^{*}$ : l'application
$\mathrm{Circ}(I) \to V^{*}$, $u \mapsto u^2$, est de degré $2$ et définit un
cube dont $\mathrm{Circ}(I)$ est l'ensemble des six faces et $V^{*}$ celui des
trois directions de faces ; $I$ est alors l'ensemble de ses quatre grandes
diagonales. C'est la réalisation classique de $\mathfrak{S}_4$ comme groupe des
rotations du cube ; le §4 la reprend en détail.

Trois ensembles se correspondent bijectivement : $V^{*}$, les sous-groupes
cycliques d'ordre $4$ et les $2$-sous-groupes de Sylow. À $x \in V^{*}$
correspondent le sous-groupe cyclique d'ordre $4$ dont $x$ est le carré du
générateur, et le $2$-Sylow $\operatorname{Cent}_G(x) = \operatorname{Norm}_G(x)$ ;
inversement, $x$ est l'élément non trivial du centre du Sylow, ou l'unique
élément d'ordre $2$ du groupe cyclique.\footnote{Le premier des deux schémas du
bas de la page 6 envoie $x \mapsto \operatorname{Cent}_G(x)$ sur les
sous-groupes cycliques d'ordre $4$ ; or le centralisateur d'une double
transposition dans $\mathfrak{S}_4$ est d'ordre $8$, c'est un $2$-Sylow. Le
second schéma, qui passe par les Sylow, a la version juste, et le triangle de la
page 19 aussi.}

\emph{Orientations.} Les signatures identifient $\mathfrak{S}_{I,\mathrm{ab}}$
et $\mathfrak{S}_{J,\mathrm{ab}}$ à $\{\pm 1\}$ ; d'où une bijection
fonctorielle $\underline{\omega}(I) \simeq \underline{\omega}(J)$, à ceci près
que les bijections compatibles aux actions de $\mathfrak{S}_I$ forment elles-mêmes
un torseur sous $\{\pm 1\}$ : il faut une convention, et il y en a deux. Pour
$i \in I$ fixé, $\underline{\omega}(I) \simeq \underline{\omega}(I \setminus \{i\})$
en complétant un ordre total de $I \setminus \{i\}$ par $i$ placé en
\emph{premier} --- le placer en dernier serait l'autre convention.

\pagerange{7}{19}
\section*{2. Le carré combinatoire et son groupe (pages 7 à 19)}

\emph{Cinq présentations.} Un carré combinatoire $Q$ peut se donner :
\begin{enumerate}
\item[(1)] par un triple $(S, A, R)$ --- quatre sommets, quatre arêtes, une
relation d'incidence ;
\item[(2)] par un ensemble $S$ à quatre éléments muni d'une involution sans
point fixe $\underline{a}$, l'\emph{antipodie}, d'où le revêtement de degré $2$
$S \to D = S/\underline{a}$ sur l'ensemble des deux \emph{diagonales} ;
\item[(3)] dualement, par $A \to C = A/\underline{a}$, où $C$ est l'ensemble des
deux \emph{codiagonales}, c'est-à-dire des paires de côtés opposés ;
\item[(2 bis), (3 bis)] par une structure de torseur sur $S$ (resp.\ sur $A$)
sous un groupe cyclique d'ordre $4$ donné comme sous-groupe de
$\mathfrak{S}_S$ (resp.\ $\mathfrak{S}_A$).
\end{enumerate}
On passe de (2 bis) à (2) en prenant pour $\underline{a}$ l'élément d'ordre $2$
du groupe cyclique ; de (2) à (1) en prenant pour arêtes les quatre parties à
deux éléments de $S$ qui ne sont pas des diagonales, avec l'incidence
tautologique ; de (3) à (1) symétriquement ; et de (1) aux autres par le
groupe $\mathfrak{D} = \mathfrak{D}_Q = \operatorname{Aut} Q$, d'ordre $8$,
isomorphe au groupe diédral $\mathbb{D}_4$.\footnote{Le début de la page 7, qui
situe ces axiomes parmi ceux des polygones réguliers, est d'une lecture très
incertaine ; les trois présentations et les deux flèches « degré 2 » sont
sûres. « cinq » (descriptions) est ajouté en interligne et n'est pas sûr.}

\emph{La structure de $\mathfrak{D}$.} Les images inverses dans $\mathfrak{D}$
des groupes de Klein $V_S \subset \mathfrak{S}_S$ et $V_A \subset \mathfrak{S}_A$
sont deux sous-groupes d'ordre $4$, encore notés $V_S$ et $V_A$. Le centre de
$\mathfrak{D}$ est $\{1, \underline{a}\}$. On a
\[
V_S \simeq \mathbb{F}_2^{C}, \qquad V_A \simeq \mathbb{F}_2^{D},
\]
l'antipodie étant dans les deux cas l'élément diagonal : les deux éléments de
$V_S \setminus \{1, \underline{a}\}$ sont les symétries par rapport aux deux
axes transverses aux paires de côtés opposés, indexées par $C$ ; ceux de
$V_A \setminus \{1, \underline{a}\}$ sont les symétries par rapport aux
diagonales, indexées par $D$. Le troisième sous-groupe d'indice $2$ est le
groupe cyclique $\mathfrak{D}^{+}$ des rotations ; ses deux générateurs forment
l'ensemble $\underline{\omega}_Q$ des orientations de $Q$, et
$\mathfrak{D}^{+} \simeq (\mathbb{Z}/4\mathbb{Z}) \wedge_{\{\pm 1\}}
\underline{\omega}_Q$. $S$ et $A$ sont des torseurs sous $\mathfrak{D}^{+}$.

L'abélianisé $\mathfrak{D}_{\mathrm{ab}} = \mathfrak{D}/\{1, \underline{a}\}$
est un plan sur $\mathbb{F}_2$, dont les trois droites sont les images de
$V_S$, $V_A$ et $\mathfrak{D}^{+}$ ; d'où un isomorphisme canonique
$\mathfrak{D}_{\mathrm{ab}} \simeq \mathbb{F}_2 \times \mathbb{F}_2$, envoyant
les images de $V_S$, $V_A$, $\mathfrak{D}^{+}$ sur la première, la seconde
droite et la diagonale.\footnote{La page écrit à deux reprises
$\mathfrak{D}^{+}_{\mathrm{ab}}$, avec un « $+$ » net ; le sens --- quotient de
$\mathfrak{D}$ par son centre, de rang $2$ --- est celui de
$\mathfrak{D}_{\mathrm{ab}}$, le groupe cyclique $\mathfrak{D}^{+}$ étant son
propre abélianisé.} Les images inverses des trois éléments non nuls sont
\[
(1,0) \leadsto C, \qquad (0,1) \leadsto D, \qquad (1,1) \leadsto
\underline{\omega}_Q,
\]
d'où la décomposition de $\mathfrak{D}$ en classes de conjugaison
$\{1\} \amalg \{\underline{a}\} \amalg C \amalg D \amalg \underline{\omega}_Q$.
Les sous-groupes propres non triviaux sont les trois sous-groupes d'indice $2$
et les cinq sous-groupes d'ordre $2$ : le centre, et deux paires contenues
respectivement dans $V_S$ et dans $V_A$. Un carré porte ainsi trois ensembles
à deux éléments remarquables, $D$, $C$ et $\underline{\omega}_Q$.\footnote{En
bas de la page 10, les deux premiers symboles de l'énumération ressemblent à
un $\mathfrak{B}$ et à un $S$ ; la suite (« en plus des deux diagonales et des
deux codiagonales ») et la page 11 montrent qu'il s'agit de $D$ et de $C$.}

\emph{Les relations entre ces trois ensembles.} D'abord
$\underline{\omega}_S \simeq C$ et $\underline{\omega}_A \simeq D$ : orienter
l'ensemble à quatre éléments $S$ revient, par la bijection
$\underline{\omega}(S) \simeq \underline{\omega}(V_S^{*})$ du §1, à orienter
l'ensemble à trois éléments $V_S^{*} = \{\underline{a}\} \amalg C$, et, l'élément
$\underline{a}$ étant distingué, à ordonner $C$, c'est-à-dire à choisir un
élément de $C$ --- moyennant, chaque fois, une convention. Ensuite, la relation
remarquable
\[
(*) \qquad \underline{\omega}_Q \wedge C \wedge D \simeq \mathbb{1},
\]
qui vient de ce que $\underline{\omega}_Q$, $C$, $D$ sont les images inverses,
dans l'extension centrale $\mathfrak{D}$ de $\mathfrak{D}_{\mathrm{ab}}$ par
$\{1, \underline{a}\}$, de trois éléments de somme nulle : le produit de trois
éléments $\varpi \in \underline{\omega}_Q$, $c \in C$, $d \in D$ tombe dans le
centre, et l'isomorphisme $(*)$ est $\varpi \wedge c \wedge d \mapsto \varpi c d
\in \{1, \underline{a}\} = \mathbb{F}_2$. En termes de groupe structural :
$(*)$ réduit à $\{x + y + z = 0\}$ le groupe $\mathbb{F}_2^3$ du torseur
$\underline{\omega}_Q \times C \times D$. Comme $c$ et $d$ ne commutent pas
($[c, d] = \underline{a}$), l'isomorphisme dépend de l'ordre des facteurs :
c'est une convention, et elle s'inverse quand on passe au carré dual (sommets
et arêtes échangés).

Les repères du carré --- couples formés d'un sommet et d'une arête incidente
--- forment un ensemble $R$ à huit éléments, le même pour $Q$ et son dual ; on a
$R/\underline{a} \simeq C \times D$, et les deux classes d'orientation de
repères, stables chacune par $\underline{a}$, s'identifient à
$\underline{\omega}_Q$. Pour vérifier que $(*)$ coïncide avec la relation
$C \wedge D \simeq \underline{\omega}_Q$ ainsi obtenue, il faut expliciter « de
bout en bout » les conventions : à une codiagonale $c$ on associe la symétrie
qui fixe chacun de ses deux côtés, à une diagonale $d$ la symétrie qui la fixe,
à un repère $(s, a)$ la rotation qui envoie $s$ sur l'autre extrémité de $a$.
Avec ces conventions, la relation $\varpi c d = 1$ dit que la classe
d'orientation $R_\varpi$ contient l'image inverse de $(c, d)$, pourvu que les
produits se lisent de gauche à droite, comme des actions à droite.\footnote{Ce
point est le nôtre ; la page dit seulement « on vérifie ». Sommets $1, 2, 3, 4$
dans l'ordre circulaire, repère $(1, \{1,2\})$ : $\varpi = (1\,2\,3\,4)$,
$c = (1\,2)(3\,4)$, $d = (2\,4)$. Le composé « $\varpi$, puis $c$, puis $d$ »
est l'identité ; le composé usuel $\varpi \circ c \circ d$ envoie $1$ sur $3$,
et vaut $\underline{a}$. Avec la composition usuelle, la relation s'écrit
$d c \varpi = 1$.} Que le passage au carré dual semble inverser la relation
vient, explique-t-il, de ce que les repères de $Q$ et de son dual sont les
mêmes, mais qu'une même classe de repères définit sur $Q$ et sur le dual deux
orientations \emph{opposées} : envoyer $s$ sur $s'$, ou $a$ sur $a'$, se font
par deux rotations de sens contraires.

\emph{Le groupe abstrait.} Comme groupe abstrait, $\mathfrak{D}$ est extension
de $W = \mathfrak{D}_{\mathrm{ab}}$ par $\mathbb{F}_2$, avec un élément non nul
distingué $\varpi \in W$ (l'image des éléments d'ordre $4$), les deux autres
éléments non nuls étant échangés par $\operatorname{Aut} \mathfrak{D}$. D'où
\[
1 \to \check{W} \to \operatorname{Aut} \mathfrak{D} \to \{\pm 1\} \to 1,
\]
où $\check{W} = \operatorname{Hom}(W, \mathbb{F}_2)$ est le groupe des
automorphismes de l'extension qui induisent l'identité sur $W$ et sur
$\mathbb{F}_2$ ; ce sont exactement les automorphismes intérieurs, et
$\check{W} \simeq W$ canoniquement (comme pour tout plan sur $\mathbb{F}_2$, par
la forme alternée $x \wedge y$). $\operatorname{Aut} \mathfrak{D}$ est produit
semi-direct de $\{\pm 1\}$ par $\check{W}$ et se trouve isomorphe à
$\mathbb{D}_4$. Le groupe $\mathfrak{D}$ seul ne permet donc de reconstruire
$Q$ qu'« à l'antipodie près » : pour le carré épinglé $Q_0$,
$1 \to \mathbb{F}_2 \to \mathbb{D}_4 = \operatorname{Aut}(Q_0) \to
\mathbb{F}_2^2 \to 1$, la donnée de $\operatorname{Aut} Q$ avec ses deux classes
de symétries marquées revient à celle d'un torseur sous le quotient
$\mathbb{F}_2^2$, c'est-à-dire du couple $(D, C)$, tandis qu'un carré est un
torseur sous $\mathbb{D}_4$ lui-même. On retrouve $Q$ à partir de
$\mathfrak{D}$ en relevant en une action transitive sur un ensemble $S$ à quatre
éléments, avec $S/\underline{a} \simeq D$, l'action de $\mathfrak{D}$ sur
$D \subset \mathfrak{D}$ ; dualement avec $C$.\footnote{La page écrit
« $D \hookrightarrow \mathfrak{D}^{+}$ » et
« $C \hookrightarrow \mathfrak{D}^{+}$ », le « $+$ » net ; or $D \subset V_A$ et
$C \subset V_S$. On lit $\mathfrak{D}$.}

\emph{Remarque} (page 15) : $\operatorname{Aut}(\mathbb{D}_4) \simeq \mathbb{D}_4$,
non par les automorphismes intérieurs, ce qui « indique » une équivalence de
groupoïdes entre carrés et groupes isomorphes à $\mathbb{D}_4$, qui n'est pas
$Q \mapsto \operatorname{Aut} Q$ et « mériterait d'être comprise ». La page 16
en donne une : à un groupe $\mathfrak{D} \simeq \mathbb{D}_4$ on associe le
carré $Q_{\mathfrak{D}}$ dont les sommets sont les quatre éléments d'ordre $2$
non centraux, l'antipodie étant $x \mapsto x\underline{a}$. Ce foncteur est
bien une équivalence de groupoïdes\footnote{Le début de la phrase de la page 16
est surchargé et illisible : on ne peut dire si la page affirme ou nie
l'équivalence. Elle a lieu (vérification nôtre) : un automorphisme de
$\mathfrak{D}$ est déterminé par son action sur ses éléments d'ordre $2$ non
centraux, qui l'engendrent ; d'où une injection
$\operatorname{Aut} \mathfrak{D} \to \operatorname{Aut} Q_{\mathfrak{D}}$ entre
deux groupes d'ordre $8$, donc une bijection.}, et la page en donne le
quasi-inverse : $\mathfrak{D}$ est engendré par $\underline{a}$ et les
sommets $x$, avec les relations $[x, y] = \underline{a}$ pour $y \neq x, x'$.
Appliqué à $\mathfrak{D} = \operatorname{Aut} Q$, ce foncteur rend un carré
dont les sommets sont $C \amalg D$ --- un autre carré que $Q$, comme le note la
marge.

\emph{Le carré par ses sommets.} Un carré équivaut à un ensemble $S$ à quatre
éléments --- un torseur sous un plan $V_S$ sur $\mathbb{F}_2$ --- et un élément
$\underline{a} \in V_S^{*}$. Pour un plan $V$, la donnée de
$\underline{a} \in V^{*}$ revient à celle de la base $V^{*} \setminus
\{\underline{a}\}$, donc d'un ensemble $C$ à deux éléments avec
$V \simeq \mathbb{F}_2^{C}$ et $\underline{a}$ diagonal. Un carré équivaut donc
à un couple $(C, S)$, $S$ torseur sous $\mathbb{F}_2^{C}$, ou encore à une
famille de $\mathbb{F}_2$-torseurs indexée par $C$, c'est-à-dire à un
revêtement de degré $2$, $A \to C$ : les sommets sont les sections,
$S \simeq \prod_{c \in C} A_c \subset \mathfrak{P}_2(A)$, d'où l'incidence.
Dualement $A \simeq \prod_{d \in D} S_d$.

\emph{Formulaire} (pages 18 et 19). Pour un carré $Q$, avec ses codiagonales
$C$, ses diagonales $D$, ses orientations $\underline{\omega}$, et les
revêtements $S \to D$, $A \to C$ :
\[
\begin{array}{ll}
(1) & \underline{\omega}_S \simeq C, \quad \underline{\omega}_A \simeq D \\
(2) & \underline{\omega} \wedge C \wedge D \simeq \mathbb{1}, \ \text{i.e.}\
\underline{\omega} \simeq C \wedge D,\ C \simeq D \wedge \underline{\omega},\
D \simeq \underline{\omega} \wedge C \\
(3) & D \simeq S/\underline{a}_S, \quad C \simeq A/\underline{a}_A \\
(4) & S \simeq \prod_{c \in C} A_c, \quad A \simeq \prod_{d \in D} S_d \\
(5) & D \simeq \bigwedge_{c \in C} A_c, \quad C \simeq \bigwedge_{d \in D} S_d
\end{array}
\]
La relation (5) suit de (3) et (4) : le produit contracté des deux
$\mathbb{F}_2$-torseurs $A_c$ est le quotient de $\prod A_c = S$ par l'élément
diagonal, qui est $\underline{a}$.

\pagerange{19}{21}
\section*{3. Les sous-groupes de $\mathfrak{S}_4$ (pages 19 à 21)}

Soit $G = \mathfrak{S}_I$, $V$ son sous-groupe de Klein distingué, de sorte que
$V^{*}$ s'identifie aux partitions de type $(2,2)$. Le triangle de la page 19
identifie entre eux $V^{*}$, les sous-groupes cycliques d'ordre $4$, les
$2$-sous-groupes de Sylow --- et les structures de carré sur $I$, dont
$\mathfrak{D}$ est alors le groupe d'automorphismes : tout groupe cyclique
d'ordre $4$ a un unique élément d'ordre $2$, qui est dans $V^{*}$ et dont il est
le centralisateur ; tout Sylow a un unique élément central non trivial, dans
$V^{*}$, dont il est le centralisateur et le normalisateur ; tout Sylow a un
unique sous-groupe cyclique d'ordre $4$, dont il est le normalisateur.

Le recensement des sous-groupes $\neq \{e\}, G$ :
\begin{enumerate}
\item[a)] les $3$ sous-groupes engendrés par un élément de $V^{*}$ ;
\item[b)] les $6$ engendrés par une transposition, groupés en trois paires
indexées par $V^{*}$ ($\tau\tau' \in V^{*}$, ou : contenus dans un même Sylow) ;
\item[c)] les $3$ sous-groupes cycliques d'ordre $4$ ;
\item[d)] $V$ lui-même ;\footnote{La page écrit $V^{*}$ ; le compte « 1 » de la
marge et le sens désignent le sous-groupe $V$.}
\item[e)] les $3$ groupes de Klein non distingués, chacun contenu dans un seul
Sylow, formant une classe de conjugaison ;
\item[f)] les $3$ sous-groupes de Sylow ;
\item[g)] les $4$ sous-groupes d'ordre $3$, indexés par $I$ ;
\item[h)] les $4$ sous-groupes d'ordre $6$, isomorphes à $\mathfrak{S}_3$ :
stabilisateurs des points de $I$, normalisateurs des $3$-Sylow ;
\item[i)] le groupe alterné, seul sous-groupe d'indice $2$, puisque
$G_{\mathrm{ab}} \simeq \mathbb{F}_2$.
\end{enumerate}
Il y a donc $19$ sous-groupes d'ordre une puissance de $2$ distincts de
$\{e\}$, en six classes de conjugaison, et $9$ d'ordre multiple de $3$
distincts de $G$, en trois classes : $28$ en tout, $30$ avec $\{e\}$ et $G$,
répartis en $11$ classes de conjugaison, qui correspondent aux classes
d'isomorphie d'espaces homogènes sous $G$.\footnote{La page écrit « 17 », ce qui
contredit son propre total : $17 + 9 = 26$, non $28$. Le compte
$3 + 6 + 3 + 1 + 3 + 3 = 19$ est celui de ses propres items a) à f), tel que
la marge le tient.} « Pour bien faire, il faudrait les relations dans un
diagramme d'inclusions » : ce diagramme n'est pas fait.\footnote{Une note
marginale du titre de la page 19 demande de « compléter par les
$3$-sous-groupes de Sylow » ; c'est ce que fait la page 21, items g) à i).}

\pagerange{22}{24}
\section*{4. Le cube et l'octaèdre attachés à un ensemble à quatre éléments (pages 22 à 24)}

L'application de degré $2$, $\mathrm{Circ}(I) \to J$, $u \mapsto u^2$, définit
un cube --- de dimension $3$ cette fois --- dont les faces sont les six
éléments de $\mathrm{Circ}(I)$, les sommets les huit sections de
$\mathrm{Circ}(I) \to J$, les arêtes les douze sections partielles au-dessus
des parties à deux éléments de $J$, l'incidence étant le prolongement des
sections ; $G$ opère sur ce cube par rotations, $G \simeq
\operatorname{Aut}^{+}(\text{cube})$.\footnote{Le paragraphe de la page 22 est
surchargé d'ajouts et de biffures ; la description du cube est celle que le sens
impose, et c'est celle que la transcription restitue dans sa note.} L'autre
application de degré $2$ vers $J$, celle des transpositions, ne définit pas de
cube : les espaces homogènes $\mathrm{Circ}(I)$ et $\mathfrak{P}_2(I)$ ne sont
pas isomorphes, leurs stabilisateurs étant $\mathbb{Z}/4\mathbb{Z}$ et
$(\mathbb{Z}/2\mathbb{Z})^2$.

Il y aurait lieu d'expliciter : les huit sommets et les huit éléments d'ordre
$3$ ; les six paires d'arêtes antipodiques et les six transpositions ; les six
faces et les six permutations circulaires ; et une \emph{orientation canonique}
du cube, c'est-à-dire un isomorphisme
$\underline{\omega}_{\text{cube}} \simeq \underline{\omega}_J$. Pour $i \in J$,
soit $Q_i$ la structure de carré sur $I$ définie par $i$ et $D_i = I/a_i$
l'ensemble de ses deux diagonales. L'injection $I \hookrightarrow \prod_{i \in J}
D_i$ fait de $I$ un torseur sous $V_J(\mathbb{F}_2) \subset \mathbb{F}_2^{J}$,
ce qui trivialise canoniquement $\bigwedge_{i \in J} D_i$ ; c'est l'ingrédient
qui permet d'écrire $\underline{\omega}_{\text{cube}} \simeq \bigwedge_{i}
\underline{\omega}(Q_i)$.\footnote{Les légendes sous les formules de la page 23
sont écrites en petit, l'une sur l'autre, et lues sous réserve ; on n'en garde
que ce que la suite (« OK », page 24) établit.}

\emph{Le compte des espaces homogènes.} Les dix classes de sous-groupes
$\neq G$ du §3 doivent correspondre à des éléments de structure du cube ; la
page donne les cardinaux $8, 12, 6, 4, 6, 3, 24$, puis un ensemble à deux
éléments, les diagonales de faces, et les couples d'éléments distincts de $J$,
et conclut « le compte y est ». En clair, avec les
stabilisateurs :\footnote{La liste est la nôtre ; la page n'est lue que par
fragments à cet endroit, et ses sept cardinaux ne sont pas légendés.}
\begin{enumerate}
\item[--] sommets ($8$), $\mathbb{Z}/3\mathbb{Z}$ ; arêtes ($12$), une
transposition ; faces ($6$), $\mathbb{Z}/4\mathbb{Z}$ ;
\item[--] grandes diagonales ($4$), $\mathfrak{S}_3$ ; paires d'arêtes opposées
($6$), un groupe de Klein non distingué ; paires de faces opposées ($3$), un
$2$-Sylow ;
\item[--] repères ($24$), $\{e\}$ ; les deux tétraèdres inscrits ($2$), le
groupe alterné ; diagonales de faces ($12$), un élément de $V^{*}$ ; couples de
directions de faces distinctes ($6$), $V$.
\end{enumerate}

\pagerange{25}{31}
\section*{5. Les $\mathbb{Z}/4\mathbb{Z}$-modules libres de rang $2$ (pages 25 à 31)}

Soit $\Lambda = \mathbb{Z}/4\mathbb{Z}$, $M$ un $\Lambda$-module libre de rang
$2$, $M_0 = M/2M$, plan sur $\mathbb{F}_2$, et $J = M_0^{*}$ l'ensemble de ses
trois éléments non nuls, notés $x_i$, $i \in J$. La multiplication par $2$
induit $M_0 \simeq 2M$, d'où l'extension de $\Lambda$-modules
$0 \to M_0 \to M \to M_0 \to 0$. Pour $x \in M_0$, la fibre $M_x$ est un torseur
sous $M_0$. Posons $M_i = M_{x_i}$. Comme $\sum_i x_i = 0$, la somme des
relèvements tombe dans $2M \simeq M_0$, et
\[
P_M = \Bigl\{ (\xi_i)_{i \in J} \in \textstyle\prod_i M_i \Bigm|
\sum_i \xi_i = 0 \Bigr\}
\]
est un torseur sous $V_J(M_0)$ ; la donnée des $M_i$ et de la trivialisation
canonique de leur somme contractée revient à celle de $P_M$.

\emph{L'algèbre $\operatorname{End}(M_0)$.} Pour tout $\mathbb{F}_2$-vectoriel
$W$, $\operatorname{Hom}(M_0, W) \simeq V_J(W)$, car
$M_0 \simeq \mathbb{F}_2^{J}/\mathrm{diag}$ ; en particulier
$\operatorname{End}(M_0) \simeq V_J(M_0)$, $u \mapsto (u(x_i))_i$, de sorte que
$P_M$ est un torseur sous $\operatorname{End}(M_0)$. Soit $\sigma_i$ la
transvection qui fixe $x_i$ et échange les deux autres ; comme
$\sum \sigma_i = 0$, $x_i \mapsto \sigma_i$ définit $M_0 \hookrightarrow
\operatorname{End}(M_0)$. Soit $\underline{\omega} = \underline{\omega}_J =
\{\rho, \rho^{-1}\}$, les deux éléments d'ordre $3$ de
$\operatorname{GL}(M_0) \simeq \mathfrak{S}_J$ ; comme $\rho + \rho^{-1} =
\mathrm{id}$, $\mathbb{F}_2^{\underline{\omega}} \hookrightarrow
\operatorname{End}(M_0)$ envoie l'élément diagonal sur l'identité. Pour
$i \neq j$, $\sigma_i, \sigma_j, \rho, \rho^{-1}$ forment une base, d'où une
décomposition en $\mathfrak{S}_J$-modules
\[
\operatorname{End}(M_0) = M_0 \oplus \mathbb{F}_2^{\underline{\omega}},
\]
les deux facteurs étant orthogonaux pour la forme $\operatorname{Tr}(uv)$ ; et
$\mathfrak{sl}(M_0) = \operatorname{Ker} \operatorname{Tr} = M_0 \oplus
\mathbb{F}_2 \cdot \mathrm{id}$, $\mathbb{F}_2^{\underline{\omega}} \cap
\mathfrak{sl}(M_0) = \mathbb{F}_2 \cdot \mathrm{id}$. Le torseur $P_M$ se
scinde donc en un couple : $I_M = P_M / \mathbb{F}_2^{\underline{\omega}}$,
torseur sous $M_0$, et $S_M = P_M / M_0$, torseur sous
$\mathbb{F}_2^{\underline{\omega}}$.

\begin{quote}
\textbf{Théorème} (page 27). Le foncteur $M \mapsto (M_0, I_M, S_M)$ est une
équivalence du groupoïde des $\Lambda$-modules libres de rang $2$ vers celui
des triplets $(M_0, I, S)$ : $M_0$ plan sur $\mathbb{F}_2$, $I$ torseur sous
$M_0$, $S$ torseur sous $\mathbb{F}_2^{\underline{\omega}(M_0)}$.
\end{quote}

\emph{Corollaire.} Les automorphismes s'en déduisent : la suite
$1 \to \operatorname{End}(M_0) \to \operatorname{GL}_\Lambda(M) \to
\operatorname{GL}(M_0) \to 1$, où $u \in \operatorname{End}(M_0)$ s'envoie sur
$\mathrm{id}_M + 2\tilde u$ pour un relèvement quelconque $\tilde u$ de
$u$,\footnote{La page dit que le composé avec
$\operatorname{End}_\Lambda(M) \to \operatorname{End}(M_0)$ est
$u \mapsto \mathrm{id}_M + 2u$, ce qu'il faut lire ainsi ; $2\tilde u$ ne
dépend pas du relèvement.} s'identifie à $1 \to M_0 \times
\mathbb{F}_2^{\underline{\omega}} \to \operatorname{Aut}(M_0, I_M, S_M) \to
\operatorname{GL}(M_0) \to 1$ : $\mathrm{id} + 2\tilde u$ agit sur $P_M$ par
translation par $(u(x_i))_i$. Les deux groupes sont d'ordre $96$. La marge
traduit : $\operatorname{GL}(2, \mathbb{Z}/4\mathbb{Z})$ est produit semi-direct
de $\operatorname{GL}(2, \mathbb{F}_2) \simeq \mathfrak{S}_3$ par
$M_2(\mathbb{F}_2) \simeq \mathbb{F}_2^2 \times
\mathbb{F}_2^{\underline{\omega}}$, et
$\operatorname{SL}(2, \mathbb{Z}/4\mathbb{Z}) = \{ g \cdot u \mid
\mathrm{sg}(g)\, e(\operatorname{Tr} u) = 1 \}$, $e : \mathbb{F}_2 \simeq
\{\pm 1\}$, puisque $\det(1 + 2u) = 1 + 2 \operatorname{Tr} u$.

\emph{Corollaire 2.} Par le §1, un couple $(M_0, I)$ est un ensemble $I$ à
quatre éléments ($M_0 \simeq V_I$) ; par le §2, un torseur sous $\mathbb{F}_2^C$
avec $C$ à deux éléments est un carré dont $C$ est l'ensemble des codiagonales.
Comme $\underline{\omega}(M_0) = \underline{\omega}(J) \simeq
\underline{\omega}(I)$, on obtient une équivalence canonique
\[
(\Lambda\text{-modules libres de rang } 2) \xrightarrow{\;\approx\;}
\{(I, Q, k)\}, \qquad k : \underline{\omega}(I) \xrightarrow{\;\sim\;} C(Q),
\]
$I$ ensemble à quatre éléments, $Q$ carré combinatoire.

\emph{Corollaire 3.} Si $M$ correspond à $(I, Q, k)$, le carré

\begin{tikzcd}
\operatorname{GL}_\Lambda(M) \arrow[r] \arrow[d] & \mathfrak{D}_Q \arrow[d, "\chi_C"] \\
\mathfrak{S}_I \arrow[r, "\mathrm{sg}"'] & \{\pm 1\}
\end{tikzcd}

est cartésien, $\chi_C$ étant l'action sur les codiagonales :
$\operatorname{GL}_\Lambda(M) \simeq \mathfrak{S}_I \times_{\{\pm 1\}}
\mathfrak{D}_Q$, et $24 \cdot 8 / 2 = 96$. Le noyau de
$\operatorname{GL}_\Lambda(M) \to \mathfrak{D}_Q$ est $\mathfrak{S}_I^{+}$, celui
de $\operatorname{GL}_\Lambda(M) \to \mathfrak{S}_I$ est
$\operatorname{Ker} \chi_C = V_S \simeq \mathbb{F}_2^{\underline{\omega}}$, et
celui du composé vers $\{\pm 1\}$ est $\mathfrak{S}_I^{+} \times
\mathbb{F}_2^{\underline{\omega}}$.

\emph{Proposition} (page 29). Le noyau $\mathfrak{S}_I^{+}$ de
$\operatorname{GL}_\Lambda(M) \to \mathfrak{D}_Q$ est aussi le groupe dérivé du
sous-groupe $\operatorname{SL}_\Lambda(M)$, d'indice $2$
($\Lambda^{*} = \{\pm 1\}$), et il est d'indice $4$ dans
$\operatorname{SL}_\Lambda(M)$, dont l'abélianisé est cyclique d'ordre $4$ ; le
déterminant est le composé de $\operatorname{GL}_\Lambda(M) \to \mathfrak{D}_Q$
et du caractère de noyau $\mathfrak{D}_Q^{+}$. Les trois caractères non
triviaux de $\mathfrak{D}_Q$ se lisent ainsi sur $g \in
\operatorname{GL}_\Lambda(M)$ d'image $\dot g$ :
\[
\det g = \chi_{\underline{\omega}}(\dot g), \qquad \mathrm{sg}\, g =
\chi_C(\dot g), \qquad \det g \cdot \mathrm{sg}\, g = \chi_D(\dot g).
\]
D'où le \emph{corollaire} (page 29) : un sous-diagramme de suites exactes pour
$\operatorname{SL}_\Lambda(M)$, de noyaux $\mathfrak{S}_I^{+}$ et
$\{1, -\mathrm{id}_M\} \simeq \{1, \underline{a}_Q\}$, et d'image
$\mathfrak{D}_Q^{+}$ ; et le \emph{corollaire 2} (page 30) :
\[
\operatorname{GL}_\Lambda(M)_{\mathrm{ab}} = \operatorname{GL}_\Lambda(M) /
(\mathfrak{S}_I^{+} \times \mathfrak{z}) \xrightarrow{\;\sim\;}
(\mathfrak{D}_Q)_{\mathrm{ab}} \simeq \mathbb{F}_2 \times \mathbb{F}_2,
\]
où $\mathfrak{z} = \{1, -\mathrm{id}_M\}$ est le centre de
$\operatorname{GL}_\Lambda(M)$ et de $\operatorname{SL}_\Lambda(M)$.

\emph{Le treillis.} La page 31 réunit l'essentiel dans le diagramme des
sous-groupes, chaque inclusion portant son indice ; les cinq carrés en sont
cartésiens :

\begin{tikzcd}[column sep=small, row sep=small, nodes={font=\scriptsize}]
 & & & \operatorname{GL}_\Lambda(M) & & \\
 & & \mathfrak{S}_I^{+} \times \mathbb{F}_2^{\underline{\omega}} \arrow[ur, "2"] & & \operatorname{SL}_\Lambda(M) \arrow[ul, "2"'] & \\
 & \mathfrak{gl}(M_0) = M_0 \times \mathbb{F}_2^{\underline{\omega}} \arrow[ur, "3"] & & \mathfrak{S}_I^{+} \times \mathfrak{z} \arrow[ul, "2"'] \arrow[ur, "2"] & & \\
 \mathbb{F}_2^{\underline{\omega}} \arrow[ur, "4"] & & \mathfrak{sl}(M_0) = M_0 \times \mathfrak{z} \arrow[ul, "2"'] \arrow[ur, "3"] & & \mathfrak{S}_I^{+} \arrow[ul, "2"'] & \\
 & \mathfrak{z} \arrow[ul, "2"'] \arrow[ur, "4"] & & M_0 \arrow[ul, "2"'] \arrow[ur, "3"] & & \\
 & & 1 \arrow[ul, "2"'] \arrow[ur, "4"] & & &
\end{tikzcd}

Les quotients que la page indique par des accolades :
$\operatorname{GL}_\Lambda(M) / \mathfrak{gl}(M_0) \simeq \mathfrak{S}_J$ ;
$\operatorname{SL}_\Lambda(M) / \mathfrak{S}_I^{+} \simeq \mathfrak{D}_Q^{+}
\simeq \operatorname{SL}_\Lambda(M)_{\mathrm{ab}}$, que la page identifie aussi à
$\bigwedge^2_\Lambda M$.\footnote{Les deux crochets de la page,
$[\mathfrak{S}_I \cdot \mathbb{F}_2^{\underline{\omega}}]$ à côté de
$\operatorname{GL}_\Lambda(M)$ et $[\mathfrak{S}_I \times \mathfrak{z}]$ à côté de
$\mathfrak{S}_I^{+} \times \mathbb{F}_2^{\underline{\omega}}$, avec une flèche
pointillée de $\mathfrak{S}_I^{+} \times \mathfrak{z}$ vers
$\operatorname{GL}_\Lambda(M)$, renvoient aux descriptions « semi-déployées » qui
suivent ; on ne les reporte pas dans le treillis. Le troisième sous-groupe
d'indice $2$, $\operatorname{Ker}(\mathrm{sg} \cdot \det)$, image inverse de
$V_A$, n'y figure pas.} Et le dictionnaire des trois ensembles à deux éléments
du carré :
\[
C_Q \simeq \underline{\omega} \leftrightarrow \mathrm{sg}, \qquad
D_Q \simeq \underline{\omega} \wedge \mu \leftrightarrow \mathrm{sg} \cdot \det,
\qquad \underline{\omega}_Q \simeq \mu \leftrightarrow \det,
\]
où $\mu = (\bigwedge^2_\Lambda M)^{*} = \operatorname{Or}(M)$ est l'ensemble des
deux générateurs de $\bigwedge^2 M$.

Les isomorphismes $\operatorname{GL}_\Lambda(M) \simeq \mathfrak{S}_I \ltimes
\mathbb{F}_2^{\underline{\omega}}$ (semi-direct) et
$\operatorname{Ker}(\mathrm{sg} \cdot \det) \simeq \mathfrak{S}_I \times
\mathbb{F}_2$ (direct) ne valent que lorsque $M$ est « semi-déployé »,
c'est-à-dire $P_M$ trivialisé : le stabilisateur d'un point de $P_M$ relève
alors $\operatorname{GL}(M_0)$, et
\[
\operatorname{GL}_\Lambda(M) \simeq \operatorname{GL}(M_0) \ltimes (M_0 \times
\mathbb{F}_2^{\underline{\omega}}) \simeq (\operatorname{GL}(M_0) \ltimes M_0)
\ltimes \mathbb{F}_2^{\underline{\omega}} \simeq \mathfrak{S}_I \ltimes
\mathbb{F}_2^{\underline{\omega}},
\]
isomorphismes qui dépendent de la trivialisation.\footnote{La page écrit le
dernier facteur en produit direct,
$(\operatorname{GL}(M_0) \cdot M_0) \times \mathbb{F}_2^{\underline{\omega}}$.
Or $\mathfrak{S}_I$ agit sur $\mathbb{F}_2^{\underline{\omega}}$ par la
signature, en échangeant $\rho$ et $\rho^{-1}$ : le produit est semi-direct,
comme le dit d'ailleurs la ligne précédente de la page.}

\pagerange{32}{39}
\section*{6. L'octaèdre d'un module, et récapitulation (pages 32 à 39)}

Soit $M^{\#} \subset M$ l'ensemble des éléments d'image non nulle dans $M_0$ ;
c'est un torseur relatif au-dessus de $J$, de groupe constant $M_0$. Le groupe
$\Lambda^{*} = \{\pm 1\}$ y opère, et
\[
\Delta_M = M^{\#}/\{\pm 1\} = \mathbb{P}(M)(\mathbb{Z}/4\mathbb{Z})
\]
est l'ensemble des six points à valeurs dans $\Lambda$ de la droite projective
définie par $M$.\footnote{La page note $M^{*}$ ou $M^{\#}$ d'un même signe ; la
transcription a choisi $M^{\#}$ pour le distinguer du groupe des unités, et on
la suit.} L'application de réduction $\Delta_M \to J =
\mathbb{P}(M)(\mathbb{F}_2)$ est de degré $2$ et fait de $\Delta_M$ les sommets
d'un \emph{octaèdre} dont $J$ est l'ensemble des trois diagonales : les douze
arêtes sont les parties à deux éléments qui ne sont pas des fibres, les huit
faces les sections de $\Delta_M \to J$. Le groupe $\operatorname{GL}(M)$ opère
par automorphismes d'octaèdre, avec pour noyau $\mathfrak{z}$, d'où
\[
\operatorname{GL}(M)' = \operatorname{GL}(M)/\mathfrak{z} = \operatorname{PGL}(M)
\xrightarrow{\;\sim\;} \operatorname{Aut}_{\mathrm{oct}}(\Delta_M) \simeq
\operatorname{Aut}^{+}(\Delta_M) \times \{1, \underline{a}_\Delta\},
\]
isomorphisme puisque les deux groupes sont d'ordre $48$. Le centre de
$\operatorname{GL}(M)'$ est $\mathbb{F}_2^{\underline{\omega}} / \mathbb{F}_2
\cdot \mathrm{id}$ : il s'envoie sur le centre trivial de $\mathfrak{S}_J$, donc
vit dans $\operatorname{End}(M_0)/\mathfrak{z} \simeq (M_0 \times
\mathbb{F}_2^{\underline{\omega}})/\mathbb{F}_2$, dont il est le sous-groupe des
invariants sous $\mathfrak{S}_J$. C'est l'antipodie $\underline{a}_\Delta$.

\emph{Les faces.} Un octaèdre de diagonales $J$ équivaut à un torseur sous
$\mathbb{F}_2^J$, l'ensemble $F(\Delta)$ de ses faces ;
$\mathbb{F}_2^J = V_J(\mathbb{F}_2) \oplus \mathbb{F}_2 \cdot \mathrm{diag}$
puisque $|J|$ est impair, l'élément diagonal agissant par l'antipodie. On a un
isomorphisme canonique
\[
P_M / \mathfrak{z} \xrightarrow{\;\sim\;} F(\Delta_M), \qquad (\xi_i)_i \longmapsto
\{\pm \xi_i\}_i ,
\]
qui redonne $\operatorname{Aut}(\Delta_M) \simeq \operatorname{GL}(M)'$ par un
isomorphisme de suites exactes : à gauche
$\operatorname{End}(M_0)/\mathbb{F}_2 \cdot \mathrm{id} \simeq \mathbb{F}_2^J
\simeq M_0 \times \mathbb{F}_2$, à droite $\mathfrak{S}_J$.

\emph{Trois ensembles à deux éléments.} À l'octaèdre sont attachés trois
ensembles à deux éléments, qui correspondent aux trois caractères non triviaux
de $\operatorname{Aut}_{\mathrm{oct}}(\Delta)_{\mathrm{ab}} \simeq \mathbb{F}_2
\times \mathbb{F}_2$ :
\begin{enumerate}
\item[--] $C_\Delta$, les deux « couleurs » des faces en damier, ou les deux
tétraèdres inscrits dans le cube dual, échangés par l'antipodie ; c'est
$\bigwedge_{i \in J} \Delta_i$, quotient de $F(\Delta)$ par $V_J(\mathbb{F}_2)$,
et, pour $\Delta = \Delta_M$, $C_\Delta = P_M / (M_0 \times \mathfrak{z})$,
de caractère $\mathrm{sg} \cdot \det$ ;
\item[--] $\underline{\omega}_\Delta = \underline{\omega}_{I_\Delta} =
\underline{\omega}_J$, où $I_\Delta = F(\Delta)/\underline{a}_\Delta$ est
l'ensemble des quatre directions de faces, de caractère $\mathrm{sg}$ ;
\item[--] $\operatorname{Or}(\Delta)$, les deux orientations de l'octaèdre, de
caractère $\det$, et $\operatorname{Or}(\Delta_M) \simeq \mu = (\det M)^{*}$.
\end{enumerate}
On a $F(\Delta) \simeq C_\Delta \times I_\Delta$ et
$\operatorname{Or}(\Delta) \wedge \underline{\omega}_\Delta \wedge C_\Delta
\simeq \mathbb{1}$. L'octaèdre \emph{orienté} associé est
$\Delta^{+} = \operatorname{Or}(\Delta) \wedge^{\{1, \underline{a}_\Delta\}}
\Delta$ ; il a $\operatorname{Or}(\Delta^{+})$ trivial,
$\underline{\omega}_{\Delta^{+}} \simeq \underline{\omega}$ et
$C_{\Delta^{+}} \simeq C_\Delta \wedge \mu \simeq \underline{\omega}$, et
l'ensemble de ses « bifaces » est $I_M = P_M / \mathbb{F}_2^{\underline{\omega}}$.%
\footnote{La page numérote (45) deux formules successives ; la suite cite (47)
et (48).}

\emph{Ce que l'octaèdre perd.} Il importe, écrit Grothendieck, de comprendre
l'extension $1 \to \mathfrak{z} \to \operatorname{GL}(M) \to \operatorname{GL}(M)'
\to 1$, et en particulier $1 \to \{\pm 1\} \to \operatorname{GL}(2,
\mathbb{Z}/4\mathbb{Z}) \to \operatorname{Aut}_{\mathrm{oct}}(\text{octaèdre
standard}) \to 1$ : c'est par sa cohomologie qu'on saura quelle donnée
supplémentaire il faut, sur un octaèdre non orienté, pour en faire l'octaèdre
d'un module $M$ --- relever un torseur sous $\operatorname{GL}(2,
\mathbb{Z}/4\mathbb{Z})'$ en un torseur sous $\operatorname{GL}(2,
\mathbb{Z}/4\mathbb{Z})$. En termes du triplet $(I, Q, k)$ : $C_\Delta \simeq
D_Q$, de sorte que l'octaèdre retient $I$ et les deux ensembles $C \simeq
\underline{\omega}_I$ et $D$, soit un torseur sous $\mathbb{D}_4 /
\{1, \underline{a}\} \simeq \mathbb{F}_2 \times \mathbb{F}_2$, qu'il faut relever
en un torseur sous $\mathbb{D}_4$ --- un carré dont on connaît déjà les
diagonales et les codiagonales.\footnote{Ce passage, entre doubles crochets sur
les pages 36 et 37, est une parenthèse que la page n'achève pas ; la lecture de
« $\mathfrak{Z}(\mathbb{D}_4)$ » est incertaine, et sa dernière phrase ne se
construit pas tout à fait.}

\emph{Récapitulation partielle} (page 38). À $M$ on associe : a) l'octaèdre
$\Delta_M = M^{\#}/\pm 1$, de diagonales $M_0^{*} = J$ ; b) l'ensemble à quatre
éléments $I_M \simeq F(\Delta)/\text{antipodie}$, avec $J \simeq
\mathfrak{P}_{2,2}(I_M)$ ; c) le carré $Q_M$ de sommets $S_M = P_M / M_0$, dont
$\underline{\omega}$ est l'ensemble des codiagonales. Les données a) et b) ne
dépendent que de la droite projective $X_M = \mathbb{P}(M)$ sur $\Lambda$ ;
$Q_M$ en dépend davantage. Et
\[
C_Q \simeq \underline{\omega}_\Delta \simeq \underline{\omega}_I \simeq
\underline{\omega}_J =: \underline{\omega}, \qquad D_Q \simeq C_\Delta \simeq
\underline{\omega} \wedge \mu, \qquad \underline{\omega}_Q \simeq
\operatorname{Or}(\Delta) \simeq (\det M)^{*} =: \mu .
\]

La page 39, d'une autre encre et prise en largeur, esquisse la version
relative au-dessus d'une base $S$, que la partie II reprendra : un revêtement
$\underline{I}$ de degré $4$ équivaut à $\underline{J}$, une section $t$ de
$\Sigma_{\underline{J}}^{*}$ et un torseur $\underline{K}$ ; $X$ et $t$,
$\underline{K}$ et $t$, sont « indépendants » ; et quand $\underline{I} \simeq
\underline{J} \amalg S$ a une section, $X \simeq \Sigma_J$, d'où un morphisme
$f : \Sigma_J \to \Sigma_J$ --- celui de la page 62.\footnote{La page n'est lue
qu'en partie, et son lien avec la récapitulation n'est pas explicite ; on la lit
d'après la suite.}

\pagerange{41}{42}
\section*{Une proposition sur $\operatorname{End}(V)$ (pages 41 et 42)}

\textbf{Proposition} (page 41). Soit $V$ un plan sur $\mathbb{F}_2$, $J = V^{*}$,
$\operatorname{Aut}(V) \simeq \mathfrak{S}_J$, $\operatorname{Aut}^{\pm}(V)$ les
éléments pairs et impairs, et $\operatorname{End}(V) \simeq M_2(\mathbb{F}_2)$.
\begin{enumerate}
\item[1)] $V' = \operatorname{Aut}^{-}(V) \cup \{0\}$ est un sous-espace de rang
$2$, canoniquement isomorphe à $V$ par $x \mapsto \tau_x$ (la transvection qui
fixe $x$).
\item[2)] $W = \operatorname{Aut}^{+}(V) \cup \{0\} = \{0, 1, \rho, \rho^{-1}\}$
est une sous-algèbre commutative de rang $2$ : c'est un corps à quatre
éléments.
\item[3)] $V'$ et $W$ sont orthogonaux pour la forme trace ; sur $V'$ elle est
l'unique forme alternée non dégénérée ($\operatorname{Tr} xy = 1$ pour $x \neq
y$ non nuls), sur $W$ elle est non alternée ($\operatorname{Tr} \rho^2 = 1$) et
non dégénérée.
\item[4)] Comme $\mathbb{F}_2[\mathfrak{S}_J]$-module, $V'$ est isomorphe à $V$
et irréductible ; $W$ est une extension non scindée $0 \to \mathbb{F}_2 \cdot 1
\to W \xrightarrow{\operatorname{Tr}} \mathbb{F}_2 \to 0$, et $\mathbb{F}_2
\cdot 1$ est son propre orthogonal dans $W$.
\item[5)] $\mathfrak{sl}(V) = \operatorname{Ker} \operatorname{Tr} = V' +
\mathbb{F}_2 \cdot 1_V$.
\end{enumerate}
C'est la proposition qui sous-tend la décomposition $\operatorname{End}(M_0) =
M_0 \oplus \mathbb{F}_2^{\underline{\omega}}$ du §5.\footnote{La page dit aussi
de $V'$ qu'il est « commutatif ». Comme sous-espace il est un groupe abélien,
mais ses éléments ne commutent pas entre eux :
$\tau_x \tau_y = \rho^{\pm 1} \neq \tau_y \tau_x$ pour $x \neq y$, et $V'$ n'est
pas stable par produit. Seul $W$ est une sous-algèbre commutative.}

\pagerange{43}{52}
\section*{II. Structures octaédrales dans $\mathbb{P}^1$ --- 1. Actions du groupe de Klein, configuration octaédrale (pages 43 à 52)}

Soit $X$ un fibré en droites projectives sur un schéma $S$ de caractéristiques
résiduelles $\neq 2$, et $V$ un groupe isomorphe à $(\mathbb{Z}/2\mathbb{Z})^2$.
On étudie les monomorphismes $V_S \hookrightarrow
\underline{\operatorname{Aut}}_S(X)$, fidèles sur chaque fibre.\footnote{La
partie II, annoncée « (25 p.) » sur sa page de titre (page 43), a sa propre
pagination à partir de la page 44. Le groupe de Klein y est noté
$\mathbb{D}_2$ au titre du §1, $\mathbb{D}_4$ aux pages 46 et 51 --- ce qui,
dans la partie I, désignait le groupe diédral d'ordre $8$.} Une involution de
$X$ est déterminée par son lieu fixe, un « axe » : diviseur étale de degré $2$.
Une base $(e_1, e_2)$ de $V$ donne deux involutions $\sigma_1, \sigma_2$ d'axes
$\Delta_1, \Delta_2$ ; elles commutent si et seulement si $\Delta_1 \cap
\Delta_2 = \emptyset$ et $\Delta_2$ est stable par $\sigma_1$, et, $\Delta_1$
étant donné, $\Delta_2$ revient à une section de $(X \setminus \Delta_1) /
\{1, \sigma_1\}$. Localement, $(X, \Delta_1) \simeq (\mathbb{P}^1_S, \{0,
\infty\})$, $\sigma_1(z) = -z$, le quotient de $\mathbb{G}_m$ par $\sigma_1$ est
$\mathbb{G}_m$ via $\lambda \mapsto \lambda^2$, et on se ramène, localement pour
la topologie étale, à la forme normale
\[
\sigma_1(z) = -z, \quad \sigma_2(z) = 1/z, \quad \sigma_3(z) = -1/z ; \qquad
\Delta_1 = \{0, \infty\}, \ \Delta_2 = \{\pm 1\}, \ \Delta_3 = \mathbb{V}(z^2 + 1).
\]
\emph{Proposition} (page 46) : localement pour la topologie étale, toute
représentation fidèle de $(\mathbb{Z}/2\mathbb{Z})^2$ dans
$\operatorname{Aut}(\mathbb{P}^1_S)$ est celle-ci.

On appelle \emph{configuration octaédrale} le diviseur étale de degré $6$
\[
\Delta = \Delta_1 \amalg \Delta_2 \amalg \Delta_3 \quad \bigl( = \{0, \infty,
\pm 1, \pm i\} \text{ dans le cas normalisé} \bigr),
\]
et plus généralement tout diviseur étale de degré $6$ localement isomorphe à
celui-ci. Sur la sphère de Riemann, ce sont les six sommets d'un octaèdre
régulier. Soit $\Gamma = \underline{\operatorname{Aut}}_S(X, \Delta)$ ; il est
net sur $S$, puisqu'il se plonge dans $\underline{\mathfrak{S}}_\Delta$. Une
structure d'octaèdre sur un ensemble à six éléments est déterminée par son
groupe d'automorphismes directs : les diagonales sont les classes de la
relation « même stabilisateur », ce stabilisateur étant cyclique d'ordre $4$ ;
les arêtes et les faces s'en déduisent. Dans le cas normalisé, le stabilisateur
de $\{0, \infty\}$ dans $\Gamma$ est $\mu_4$, et de même pour $\{\pm 1\}$,
$\{\pm i\}$ ; ces trois groupes cycliques engendrent le groupe des rotations de
l'octaèdre, d'où $\operatorname{Aut}^{+}_{\mathrm{oct}}(\Delta) \subset \Gamma$.

Pour l'inclusion inverse, la page calcule le fixateur de deux sommets non
opposés, $\infty$ et $-1$ : un tel automorphisme est $z \mapsto az + (a - 1)$,
et s'il n'est pas l'identité il permute sans point fixe $\{0, 1, i, -i\}$. En
faisant $z = 0, 1, i, -i$, on trouve que $a$ doit appartenir à chacun des
ensembles $\{2, 1 \pm i\}$, $\{\tfrac12, \tfrac{1 \pm i}{2}\}$,
$\{\tfrac{1-i}{2}, 1 - i, -i\}$, $\{\tfrac{1+i}{2}, 1 + i, i\}$, dont
l'intersection est vide si la caractéristique est $\neq 2, 5$. Alors le
fixateur est trivial, $\Gamma \hookrightarrow
\operatorname{Aut}_{\mathrm{oct}}(\Delta)$, et même $\Gamma =
\operatorname{Aut}^{+}_{\mathrm{oct}}(\Delta)$ : $\Gamma$ est étale,
localement isomorphe à $(\mathfrak{S}_4)_S$.\footnote{Les deux notes marginales
de la page 49 bornent le calcul : « $a^4 = \pm 4$, qui n'est pas égal à $1$ si
car.\ $\neq 3$ et $5$ ». En caractéristique $3$, le cas $z = i$ élimine la seule
coïncidence ($a = 2 = \tfrac12$). En caractéristique $5$ en revanche, $i \in
\mathbb{F}_5$, la configuration $\{0, \infty, \pm 1, \pm i\}$ est
$\mathbb{P}^1(\mathbb{F}_5)$ tout entier, les quatre ensembles ont en commun
$\{2, 3, 4\}$, et $\underline{\operatorname{Aut}}(X, \Delta) =
\operatorname{PGL}_2(\mathbb{F}_5) \simeq \mathfrak{S}_5$, d'ordre $120$. La
page ne revient pas sur ce cas ; l'hypothèse « caractéristiques $\neq 5$ » est
ajoutée par nous partout où elle est nécessaire.}

Une structure octaédrale sur un diviseur étale $\Delta$ de degré $6$ est la
donnée d'une antipodie $\underline{a}_\Delta$, involution sans point fixe sur
chaque fibre géométrique ; elle est \emph{régulière} si
$\underline{\operatorname{Aut}}^{+}_{\mathrm{oct}}(\Delta)$ opère sur $X$ de
façon compatible à son action sur $\Delta$ (l'action est alors unique). Les
structures octaédrales régulières sont exactement celles des configurations
octaédrales, et, $\Delta$ régulière, le sous-groupe de Klein distingué de
$\operatorname{Aut}^{+}_{\mathrm{oct}}(\Delta)$ opère fidèlement sur $X$ et
redonne $\Delta$. Ces structures sont « directes » : localement, les
automorphismes qui fixent $\{0, \infty\}$ sont les rotations par $\mu_4$. Et le
revêtement des orientations d'une configuration régulière est canoniquement
$\mu_4^{*} = \operatorname{Spec} \mathcal{O}_S[i]/(i^2 + 1)$, le schéma des
racines primitives quatrièmes de l'unité.

\begin{quote}
\textbf{Théorème} (page 51). Soit $X$ un fibré en droites projectives sur un
schéma $S$ de caractéristiques résiduelles $\neq 2$, et $G =
\underline{\operatorname{Aut}}_S(X)$, forme de $\operatorname{PGL}(2)_S$. Soient
a) l'ensemble des sous-schémas en groupes de $G$ localement isomorphes à
$(\mathbb{Z}/2\mathbb{Z})^2_S$ ; b) celui des sous-schémas en groupes
localement isomorphes à $(\mathfrak{S}_4)_S$ ; c) celui des configurations
octaédrales régulières $(\Delta, \underline{a}_\Delta)$ dans $X$. On a des
bijections canoniques compatibles :
a)~$\to$~b), $\Gamma = \underline{\operatorname{Norm}}_G \underline{V}$ ;
b)~$\to$~a), $\underline{V}$ l'unique sous-groupe distingué fini étale de degré
$4$ de $\Gamma$ ; a)~$\to$~c), $\Delta$ la somme des axes des sections de
$\underline{V}^{*}$ ; c)~$\to$~b), $\Gamma =
\underline{\operatorname{Aut}}^{+}_{\mathrm{oct}}(\Delta)$. Si de plus $5$ est
inversible sur $S$, $\underline{\operatorname{Aut}}^{+}_{\mathrm{oct}}(\Delta)
\xrightarrow{\;\sim\;} \underline{\operatorname{Aut}}_S(X, \Delta)$.
\footnote{La page écrit « configurations icosaédrales » en c), là où elle vient
de corriger « icosaédrale » en « octaédrale » en tête de page ; de même au
corollaire 1 et page 59. On lit octaédrale partout. Elle écrit aussi c)~$\to$~b)
comme $\Gamma = \underline{\operatorname{Aut}}_S(X, \Delta) \leftarrow
\underline{\operatorname{Aut}}^{+}_{S\text{-oct}}(\Delta)$, isomorphisme qui
n'a lieu que hors de la caractéristique $5$.}
\end{quote}

En marge :
$\Gamma / \underline{V} \simeq \underline{\operatorname{Aut}}_{\text{gr}}(\underline{V})
\simeq \underline{\operatorname{Aut}}_S(\underline{V}^{*})$.

\emph{Corollaire 1.} Si $5$ est inversible sur $S$, un diviseur étale de degré
$6$ porte au plus une structure octaédrale régulière, et le schéma en groupes
de ses automorphismes directs est alors $\underline{\operatorname{Aut}}_S(X,
\Delta)$.\footnote{L'hypothèse est nécessaire : sur $\mathbb{F}_5$, le diviseur
$\mathbb{P}^1(\mathbb{F}_5)$ porte cinq structures octaédrales régulières, une
pour chacun des cinq sous-groupes $\mathfrak{S}_4$ de
$\operatorname{PGL}_2(\mathbb{F}_5) \simeq \mathfrak{S}_5$.}

\emph{Corollaire 2.} La catégorie des couples $(X, \Delta)$ sur $S$, $\Delta$
configuration octaédrale régulière --- ou, ce qui revient au même, des couples
$(X, \underline{V})$, $(X, \Gamma)$ --- est équivalente à celle des octaèdres
sur $S$ munis d'un isomorphisme $\underline{\operatorname{Or}}_{\Delta/S}
\simeq \mu^{*}_{4,S}$ (une « $\mu_4^{*}$-orientation »), ou encore à celle des
revêtements finis étales de degré $4$ de $S$ : on associe à $(X, \Delta)$ le
revêtement des faces modulo antipodie, ou celui des $3$-Sylow de
$\underline{\operatorname{Aut}}^{+}_{\mathrm{oct}}(\Delta)$. La structure type
est la conique $x^2 + y^2 + z^2 = 0$ de $\mathbb{P}^2_{\mathbb{Z}[1/2]}$, avec
l'action de $\{\pm 1\}^3 / \mathrm{diag}$ par changements de signe ; elle n'a pas
de point réel, et sa biface distinguée est formée des faces $\{(0,1,i), (i,0,1),
(1,i,0)\}$ et de leurs conjuguées.

\pagerange{53}{56}
\section*{2. Courbes relatives de type $(0,4)$ (pages 53 à 56)}

Soit $\underline{I} \subset X$ un diviseur relatif étale de degré $4$ : une
droite projective à quatre points marqués non ordonnés. Le schéma en groupes
$\underline{\mathfrak{S}}_{\underline{I}}$ est localement isomorphe à
$(\mathfrak{S}_4)_S$ ; son sous-groupe de Klein distingué
$\underline{V}_{\underline{I}}$ fait de $\underline{I}$ un torseur.

\emph{Proposition} (page 53). L'action de $\underline{V}_{\underline{I}}$ sur
$\underline{I}$ se prolonge, de façon unique, en une action sur $X$.

La donnée de $\underline{I}$ équivaut à celle de $\underline{V}$ --- ou du
revêtement étale cubique $\underline{J} = \underline{V}^{*}$ --- et d'un torseur
sous $\underline{V}$ ; plonger $\underline{I}$ dans $X$ revient à faire opérer
fidèlement $\underline{V}$ sur $X$ (le §1), puis à donner une section de
$(X \setminus \Delta)/\underline{V}$, dont $\underline{I}$ est l'image inverse.
Or une structure octaédrale équivaut à :
\begin{enumerate}
\item[a)] $\underline{V}$, ou $\underline{J}$, ou encore un fibré en droites
projectives $\Sigma_{\underline{J}} = \mathbb{P}(V_{\underline{J}})$ muni du
diviseur étale $D_{\underline{J}} \simeq \underline{J}$ de degré $3$ ;
\item[b)] un torseur $\underline{K}$ sous $\underline{V}$ --- un autre
revêtement de degré $4$ que $\underline{I}$ ---, ou encore un revêtement
principal $X^{*}$ de $\Sigma_{\underline{J}}^{*} = \Sigma_{\underline{J}}
\setminus D_{\underline{J}}$ de groupe $\underline{V}$, « admissible »,
c'est-à-dire dont la classe dans $H^0(S, \mathcal{H}^1(\Sigma^{*}_{\underline{J}}
/ S, \underline{V})) \simeq \operatorname{Hom}(\underline{V}_{\underline{J}},
\underline{V})$ est l'identité, par la formule $X^{*} = \underline{K}
\wedge_{\underline{V}} \Sigma^{*}_{\underline{J}, a}$.
\end{enumerate}
Ici $\Sigma_{\underline{J}, a}$ est la conique $\sum_{i \in \underline{J}}
x_i^2 = 0$, sur laquelle $V_{\underline{J}} = \mathbb{F}_2^{\underline{J}} /
\mathbb{F}_2$ opère par changements de signe, et dont le quotient est la droite
$\sum y_i = 0$ par $y_i = x_i^2$ : un revêtement galoisien de groupe
$(\mathbb{Z}/2\mathbb{Z})^2$ de $\mathbb{P}^1$ par $\mathbb{P}^1$, ramifié au-dessus
des trois points de $D_{\underline{J}}$ --- le revêtement de Klein de la droite
projective. Ainsi
\[
(X, \underline{I}) \longleftrightarrow \bigl(\underline{J},\ \underline{K},\
x \in \Gamma(\Sigma^{*}_{\underline{J}} / S)\bigr), \qquad \underline{I} \simeq
\underline{K} \wedge_{\underline{V}} (\Sigma^{*}_{\underline{J}, a} / x),
\]
et quand $\underline{I} = I_S$ est constant et numéroté, $\underline{J} = J_S$,
et la donnée de $(X, I_S \hookrightarrow X)$ se réduit à la section $x$ de
$\Sigma^{*}_J$ --- le point de $\mathbb{P}^1$ privé de trois points qui classe
quatre points numérotés de la droite ---, $\underline{K}$ étant alors déterminé
par $X/x \simeq I_S$, c'est-à-dire $\underline{K} \simeq I_S \wedge_{\underline{V}}
(\Sigma^{*}_{J,a}/x)$.\footnote{La page 55 n'est lue qu'en partie, surtout vers
le bas, et ses deux notes marginales très incomplètement ; la relation (14)
entre $\underline{I}$, $\underline{K}$ et $x$ est dans la seconde. La première
décrit la conique comme $Q = 0$ dans $\check{\mathbb{P}}(p_*
\mathcal{O}_{\underline{J}})$, $Q(u) = \operatorname{Tr} u^2$, et l'hyperplan
comme $\operatorname{Tr} u = 0$.}

\pagerange{57}{64}
\section*{3. Relation avec les courbes elliptiques ; l'application $f^t_2$ (pages 57 à 64)}

Soit $E_0$ une courbe elliptique sur $S$ (caractéristiques résiduelles
$\neq 2$), $E$ un torseur sous $E_0$, et $\sigma$ un automorphisme de $E$
induisant $-\mathrm{id}$ sur $E_0 \simeq \underline{\operatorname{Pic}}^0_{E/S}$.
Localement, $\sigma x = a - x$, où $a$ est intrinsèquement une section de
$E \wedge_{E_0} E \simeq \underline{\operatorname{Pic}}^2_{E/S}$, et $\sigma^2 =
\mathrm{id}$. Le lieu fixe $\underline{I} = E^{\sigma} = \{x \mid 2x = a\}$ est
un torseur sous $\underline{V} = {}_2E_0$, localement isomorphe à
$(\mathbb{Z}/2\mathbb{Z})^2_S$. Le couple $(E, \sigma)$ équivaut à $(E, a)$ et à
$(E_0, \underline{I})$, avec $E = \underline{I} \wedge_{\underline{V}} E_0$ ;
l'action de $\underline{V}$ sur $E$ commute à $\sigma$, et $x \mapsto 2x - a$
identifie $E/\underline{V}$ à $E_0$ en transformant $\sigma$ en $x \mapsto -x$.

Les quotients $X = E/\{1, \sigma\}$ et $\Sigma = E_0 / \{\pm 1\}$ sont des
fibrés en droites projectives, et $X/\underline{V} \simeq \Sigma$.\footnote{Une
insertion peu nette, sous « fibrés en droites », se lit « canoniquement
isomorphes » ; son rattachement est incertain. $X$ et $\Sigma$ ne sont pas
canoniquement isomorphes en général : ils le deviennent quand $\underline{I}$
est trivialisé.} Le revêtement quadratique $E \to X$ est ramifié exactement en
$\underline{I}$, qui se plonge dans $X$ de façon compatible à $\underline{V}$ et
tombe dans l'ouvert $X^{*}$ où $\underline{V}$ opère librement. On est donc dans
la situation du §2 : $\Sigma \simeq \Sigma_{\underline{J}}$, $\underline{J} =
\underline{V}^{*}$, $\underline{I}$ est l'image inverse d'une section de
$\Sigma^{*}_{\underline{J}}$, et l'image inverse $\Delta$ de $D_{\underline{J}}$
dans $X$ est la configuration octaédrale de l'action de $\underline{V}$ sur $X$.

Ensemblistement, $\Delta$ est formé des images des points $x$ pour lesquels il
existe $\xi \in {}_2E_0$, $\xi \neq 0$, avec $x + \xi = a - x$ ; donc $4x = 2a$
mais $2x \neq a$ :
\[
\Delta = \Bigl( \underline{I} \wedge_{\underline{M}_0} \underline{M} \setminus
\underline{I} \Bigr) \Big/ \{1, \sigma\}, \qquad \underline{M} = {}_4E_0, \quad
\underline{M}_0 = 2\underline{M} = {}_2E_0 = \underline{V},
\]
de degrés $16 - 4 = 12$, puis $6$. Pour $E = E_0$ : \emph{les douze points
d'ordre exactement $4$ de $E_0$, groupés en six paires $\pm x$, forment la
configuration octaédrale de $E_0/\{\pm 1\}$.} C'est la réalisation géométrique
de l'octaèdre $\Delta_M = M^{\#}/\{\pm 1\}$ de la partie I, §6. On a le
diagramme, cartésien au-dessus de $\Sigma^{*}_{\underline{J}}$ :\footnote{La
page le dit « presque cartésien », et précise les ramifications : $E \to X$
au-dessus de $\underline{I}$, $E_0 \to \Sigma_{\underline{J}}$ au-dessus de
$D_{\underline{J}}$ et de $t$, $X \to \Sigma_{\underline{J}}$ au-dessus de
$D_{\underline{J}}$. Le carré ne peut être cartésien au-dessus de
$D_{\underline{J}}$, où $E_0 / \Sigma_{\underline{J}}$ est ramifié et $E/X$ ne
l'est pas (début de la page 61).}

\begin{tikzcd}
E \arrow[r, "\text{étale}"] \arrow[d] & E_0 \simeq E/\underline{V} \arrow[d] \\
X \arrow[r] & \Sigma_{\underline{J}}
\end{tikzcd}


\emph{L'application $f^t_2$.} Dans le cas $E = E_0$, $\sigma = -1$,
$\underline{I} = {}_2E_0 = \underline{J} \amalg t(S)$, où $t$ est l'image de
l'origine, on a $X = \Sigma_{\underline{J}}$, et le passage au quotient de
$2\,\mathrm{id}_{E_0}$ donne
\[
f^t_2 : \Sigma_{\underline{J}} \longrightarrow \Sigma_{\underline{J}},
\]
qui ne doit dépendre que de $\underline{J}$ et de $t$ --- lesquels déterminent
$E_0$ à isomorphisme non unique près, localement pour la topologie
étale.\footnote{La page le souhaite (« qui doit pouvoir se définir ») sans
donner de définition intrinsèque ; elle passe aux coordonnées. Une note
marginale : « il faudrait revenir aux descriptions intrinsèques ».} Elle est
caractérisée, pour $\underline{J} = J_S$, par : si $\sigma_i^t$ est
l'involution qui échange $t$ et $i$, et les deux autres points de $J$, et
$\Delta^t_i$ son lieu fixe, alors $f^t_2(\Delta^t_i) = \xi_i$, le point $i$ de
$D_J$ ; et $f^t_2(J \cup t) = t$. ($\sigma^t_i$ est induite par la translation
par le point d'ordre $2$ d'indice $i$, et $\Delta^t_i$ est l'image des $x$ tels
que $2x$ soit ce point.)

Pour $J = \{0, 1, \infty\} \subset \mathbb{P}^1$ --- la normalisation de
Legendre --- on trouve
\[
\sigma^t_0(z) = \frac{z - t}{z - 1}, \qquad \sigma^t_1(z) = \frac{t}{z}, \qquad
\sigma^t_\infty(z) = t\, \frac{z - 1}{z - t} = \sigma^t_0 \sigma^t_1,
\]
de lieux fixes
\[
\Delta^t_0 : z^2 - 2z + t = 0, \ z = 1 \pm \sqrt{1 - t} ; \qquad \Delta^t_1 :
z^2 = t ; \qquad \Delta^t_\infty : z^2 - 2tz + t = 0, \ z = t \pm \sqrt{t(t - 1)}.
\]
\footnote{La page écrit $\Delta^t_\infty = \{t \pm \sqrt{t(1 - t)}\}$ ; le
discriminant réduit de $z^2 - 2tz + t$ est $t^2 - t = t(t - 1)$.} Comme $f^t_2$
est de degré $4$ et envoie $\Delta^t_0, \Delta^t_1, \Delta^t_\infty$ sur $0, 1,
\infty$ avec multiplicité $2$, $f^t_2 = c_t \bigl(\frac{z^2 - 2z + t}{z^2 - 2tz
+ t}\bigr)^2$ ; en un point $\zeta$ de $\Delta^t_1$, $\zeta^2 = t$, le quotient
vaut $\frac{t - \zeta}{t(1 - \zeta)} = -\zeta/t$, d'où $c_t = t$ et
\[
f^t_2(z) = t \left( \frac{z^2 - 2z + t}{z^2 - 2tz + t} \right)^2 .
\]
On vérifie que $0, 1, \infty, t$ vont sur $t$. Les itérées $f^t_{2^n} =
(f^t_2)^{\circ n}$ correspondent à $2^n\,\mathrm{id}_{E_0}$, et l'image inverse
de $t$ par $f^t_{2^n}$ est l'image de ${}_{2^n}E_0$.

\pagerange{65}{68}
\section*{4. « Détermination de $M_{1,1}(4)$ » (pages 65 à 68)}

Soit $E$ une courbe elliptique sur $S$ (caractéristiques résiduelles $\neq 2$)
et $M = {}_4E$, faisceau étale de $\mathbb{Z}/4\mathbb{Z}$-modules localement
libres de rang $2$, $M_0 = M \otimes \mathbb{F}_2 \simeq {}_2E$. Par la partie
I, §5, devenue relative, $M$ équivaut à $M_0$ --- c'est-à-dire au revêtement
cubique $\underline{J} \simeq M_0^{*}$, $M_0 = \mathbb{F}_2^{\underline{J}} /
\mathbb{F}_2$ --- et à un torseur $P$ sous $V_{\underline{J}}(M_0) \simeq
\underline{\operatorname{End}}(M_0) \simeq M_0 \oplus
\mathbb{F}_2^{\underline{\omega}}$, où $\underline{\omega}$ est le revêtement
quadratique des orientations de $\underline{J}$. Intrinsèquement :
$V_{\underline{J}}$ est exact, d'où $0 \to V_{\underline{J}}(M_0) \to
V_{\underline{J}}(M) \to V_{\underline{J}}(M_0) \to 0$, et $P$ est l'image
inverse de la section $\mathrm{id}_{M_0}$, qui correspond à l'inclusion
$\underline{J} \hookrightarrow M_0$. Une structure de niveau $4$,
$(\mathbb{Z}/4\mathbb{Z})^2_S \simeq {}_4E$, équivaut à une trivialisation de
$\underline{J}$ et une section de $P$.

Indépendamment de $\underline{J}$, un épinglage « relatif » est une section de
$P_M$, c'est-à-dire un morphisme $\underline{J} \to M$ relevant l'inclusion
$\underline{J} \hookrightarrow M_0$ et de somme nulle : sur chaque fibre
géométrique, des points $\xi_i$ avec $2\xi_i = \varepsilon_i$ et $\sum \xi_i =
0$. Il tombe dans $M \setminus M_0$ ; composé avec $M \subset E \to
\Sigma_{\underline{J}} = E/\{\pm 1\}$, il donne $\underline{J} \to \Delta_t$,
dont le composé avec $f_t : \Delta_t \to \underline{J}$ (étale de degré $2$,
induit par la multiplication par $2$) est l'identité : c'est une \emph{face} de
l'octaèdre $\Delta_t$. L'octaèdre devient ainsi $\underline{J}$-épinglé,
canoniquement isomorphe à l'octaèdre standard défini par $\underline{J}$ et une
section de $\mu_4^{*}$, « en prenant les revêtements biquadratiques » --- et la
page s'arrête là.\footnote{Le reste de la page 68 est blanc. Le titre annonce une
détermination de l'espace de modules des courbes elliptiques à structure de
niveau $4$ ; ce qui est fait est la description de ses données, non celle de
l'espace. On sait par ailleurs (remarque nôtre) que la courbe modulaire de
niveau $4$ est la courbe « octaédrale » de Klein, de groupe
$\operatorname{PSL}_2(\mathbb{Z}/4\mathbb{Z}) \simeq \mathfrak{S}_4$ ; la page ne
le dit pas.}

\pagerange{69}{73}
\section*{III. Épinglages de Legendre --- Digression : revêtements de Kummer épinglés par les espaces tangents (pages 69 à 73)}

Soit $f : X \to S$ lisse de dimension relative $1$, $t$ une section, $X^{*} = X
\setminus t(S)$. On considère des $X' \to X$ finis, lisses sur $S$, munis d'une
section $t'$ au-dessus de $t$, et modérément ramifiés en $t'$ d'indice $n$ (avec
$n$ premier aux caractéristiques résiduelles) : localement pour la topologie
étale, $X'$ est le revêtement de Kummer $z^n = \varphi$, $\varphi$ équation
locale de $t(S)$. La norme locale donne alors un isomorphisme canonique
\[
\kappa_{t'} : T_{X'/S, t'}^{\otimes n} \xrightarrow{\;\sim\;} T_{X/S, t} .
\]
Si de plus un schéma en groupes fini étale $\underline{G}$ fait de $X'|_{X^{*}}$
un torseur sous $\underline{G}_{X^{*}}$, le stabilisateur de $t'$ opère sur
$T_{X'/S, t'}$ par $\mu_{n,S}$, d'où un monomorphisme canonique $i : \mu_{n,S}
\hookrightarrow \underline{G}$, et le torseur définit une classe $\xi \in
H^0(S, R^1 f^{\circ}_{*} \underline{G})$, $f^{\circ} = f|_{X^{*}}$.

\emph{Proposition} (page 71) : $(X', t')$, avec l'action de $\underline{G}$, est
déterminé à isomorphisme unique près par $i$, $\xi$ et $\kappa_{t'}$.
Précisément, Grothendieck énonce --- « je dis que ce foncteur est une
équivalence », sans démonstration --- que la catégorie des $(X', t',
\underline{G})$ est équivalente à celle des systèmes $(\underline{G}, n, i, \xi,
T, k)$, $T$ module inversible sur $S$ et $k : T^{\otimes n} \simeq T_{X/S,t}$.
Le cas utile est $\underline{G} = \mu_2 = \{\pm 1\}_S$, $n = 2$ : pour $\xi$
fixé, les revêtements quadratiques de $X$, étales sur $X^{*}$ et totalement
ramifiés le long de $t$, équivalent aux couples $(T, k : T^{\otimes 2} \simeq
T_{X,t})$.\footnote{Rigidifier un revêtement ramifié par une racine $n$-ième du
vecteur tangent au point de ramification est proche de ce qu'on appellera plus
tard un \emph{point base tangentiel} (Deligne, 1989) ; le rapprochement est le
nôtre.}

Pour les courbes elliptiques, $X = \Sigma^{*}_{\underline{J}}$, $X' = E$ privée
de ses points d'ordre $2$, et $t$ l'image de l'origine.\footnote{La page écrit
bien $\Sigma^{*}_{\underline{J}}$, avec l'astérisque : $E \to \Sigma_{\underline{J}}$
est ramifié en quatre points, et c'est au-dessus de $\Sigma_{\underline{J}}
\setminus D_{\underline{J}}$ qu'il n'est ramifié qu'en $t$. Les deux dernières
lignes de la page 73 ne sont lues qu'en partie.}

\pagerange{74}{79}
\section*{L'épinglage de Legendre sur un corps (pages 74 à 79)}

Soit $E$ une courbe elliptique sur un corps algébriquement clos $k$ de
caractéristique $\neq 2$ --- seule cette dernière hypothèse compte, et tout vaut
sur une base. Soit $J = {}_2E(k)^{*}$, $\Sigma_J = E/\{\pm 1\} = \mathbb{P}(V_J)$,
$V_J = V_J(k)$, la droite projective où les trois points d'ordre $2$ marquent
le diviseur $D_J$. Par la digression, la donnée de $E$ équivaut à celle de $J$,
du point $t \in \Sigma^{*}_J$ image de l'origine, d'une droite $T$ (c'est
$T_{E,0}$) et d'un isomorphisme $T^{\otimes 2} \simeq T_{\Sigma, t}$ --- comme
équivalence de groupoïdes, sur toute base de caractéristiques résiduelles
$\neq 2$.

On a canoniquement $T_\Sigma \simeq \mathcal{O}_\Sigma(2)(\underline{\omega})$,
$\underline{\omega} = \underline{\omega}(J)$ l'ensemble des deux ordres
circulaires de $J$, et $L(\underline{\omega}) = L \wedge_{\{\pm 1\}}
\underline{\omega}$ : c'est que $\bigwedge^2 V_J$ est la droite tordue par la
signature, engendrée, pour un ordre circulaire $(i, j, k)$, par $(e_i - e_j)
\wedge (e_j - e_k)$.\footnote{La convention pour $\mathbb{P}(V_J)$ est celle
qui se dégage des pages 95 à 120 : un point $t$ correspond à une droite $F_t
\subset V_J$ ($\xi_i = k\,(e_j - e_k)$), $\mathcal{O}_t(1) = V_J / F_t$, et
$F_t \simeq \mathcal{O}_t(-1)(\underline{\omega})$. L'isomorphisme $T_\Sigma \simeq
\mathcal{O}(2)(\underline{\omega})$ ne dépend pas de cette convention, la
torsion par $\underline{\omega}$ étant d'ordre $2$.}

Le candidat naturel pour $T$ est $T_0 = \mathcal{O}_t(1)$ ; mais il y a
\emph{deux} isomorphismes naturels $T_0^{\otimes 2} = \mathcal{O}_t(2) \simeq
\mathcal{O}_t(2)(\underline{\omega})$, un pour chaque élément de
$\underline{\omega}$. Tordre $T_0$ par un torseur sous $\mu_2$ n'y change rien,
puisque son carré est inchangé. On le tord donc par un torseur $Q_0$ sous
$\mu_4$ --- les quatre sommets d'un \emph{carré}, dont $\mu_4$ est le groupe des
rotations :
\[
T_0 = \mathcal{O}_t(1) \wedge_{\mu_4} Q_0, \qquad T_0^{\otimes 2} \simeq
\mathcal{O}_t(2) \wedge_{\mu_2} D_{Q_0}, \qquad D_{Q_0} = Q_0 / \{\pm 1\},
\]
$D_{Q_0}$ étant l'ensemble des deux diagonales du carré.\footnote{La page écrit
« l'un des deux diagonales » ; il s'agit de l'ensemble des deux.} Si l'on se
donne $\underline{\omega}_{Q_0} \simeq \mu_4^{*} = \{i, -i\}$ et $D_{Q_0} \simeq
\underline{\omega}$, on obtient un isomorphisme canonique $T_0^{\otimes 2}
\simeq T_{\Sigma,t}$.

\emph{Définition.} Un \emph{épinglage de Legendre} de $E$, subordonné au carré
$Q_0$, est un isomorphisme $\kappa : T_0 \xrightarrow{\;\sim\;} T_{E,0}$ tel que
$\kappa^{\otimes 2}$ soit compatible aux deux isomorphismes vers $T_{\Sigma,t}$
:

\begin{tikzcd}
T_0^{\otimes 2} \arrow[rr, "\kappa^{\otimes 2}"] \arrow[dr] & & T_{E,0}^{\otimes 2} \arrow[dl] \\
 & T_{\Sigma,t} &
\end{tikzcd}

Il y en a exactement deux, opposés. Pour $E$ variable, $\underline{\omega} =
\underline{\omega}_J$ n'est pas fixé, et un épinglage de Legendre est un couple
$(\alpha, \kappa)$, $\alpha : \underline{\omega}_J \simeq D_{Q_0}$, $\kappa$ comme
ci-dessus.

Le carré $Q_0$ est choisi une fois pour toutes, avec pour seule donnée
$\underline{\omega}_{Q_0} \simeq \mu_4^{*}$, c'est-à-dire $\mu_4^{*} \simeq
\operatorname{Aut}^{+}(Q_0)$ ; le choix le plus commode est un carré dont un
côté a ses deux extrémités indexées par $\mu_4^{*} = \{\iota, \iota'\}$, avec
$s_{\iota'} = \iota\, s_\iota$ (d'où $s_\iota = \iota' s_{\iota'}$, puisque
$\iota' = \iota^{-1} = -\iota$). Alors $D_{Q_0} \simeq \mu_4^{*}$, et la donnée
$\alpha$ devient un isomorphisme $\underline{\omega}_J \simeq \mu_4^{*}$.%
\footnote{On note ici $\iota, \iota'$ les deux racines carrées de $-1$, que la
page note $i, i'$, pour réserver $i$ aux éléments de $J$. Un tel $Q_0$ existe
sur $\mathbb{Z}[1/2]$ (remarque nôtre) : le schéma des racines de $x^4 = -4$,
c'est-à-dire des sommets $\pm 1 \pm i$ du carré, avec $s_\iota = 1 - \iota$ ;
on a bien $\iota(1 - \iota) = 1 + \iota = s_{\iota'}$, et $x \mapsto -x^2/2$
identifie $Q_0/\{\pm 1\}$ à $\mu_4^{*}$ en envoyant $s_\iota$ sur $\iota$.}

\begin{quote}
La catégorie (un groupoïde) des courbes elliptiques sur $k$ munies d'un
épinglage de Legendre est équivalente à celle des triplets $(J, t, \alpha)$ :
$J$ à trois éléments, $t \in \Sigma^{*}_J(k)$, $\alpha : \underline{\omega}_J
\simeq \mu_4^{*}$.
\end{quote}

Un \emph{épinglage de Legendre--Jacobi d'échelon $2$} ajoute un épinglage de
Jacobi d'échelon $2$ --- une bijection $J \simeq \{0, 1, \infty\}$, qui identifie
$\Sigma_J$ à $\mathbb{P}^1_k$. Alors $\underline{\omega}_J$ a un élément
privilégié, et $\alpha$ revient au choix d'une racine carrée $i$ de $-1$ : c'est
« la donnée arithmétique ». La catégorie des courbes munies d'un tel épinglage
est rigide, et équivalente à celle des couples $(t, i)$, $t \in U_{0,3}(k) =
\mathbb{P}^1(k) \setminus \{0, 1, \infty\}$, $i \in \mu_4^{*}$. C'est la famille
de Legendre, $\lambda = t$, dont le retournement est tué au prix d'un choix de
$i$.

\pagerange{79}{92}
\section*{L'épinglage de Legendre sur une base ; le schéma de modules (pages 79 à 92)}

On travaille au-dessus de $S_0 = \operatorname{Spec} \mathbb{Z}[1/2]$, où $Q_0$
est un $\mu_4$-torseur. Pour un $\mathcal{O}_S$-module $L$ on pose
\[
L^{\natural} = L \wedge_{\mu_{4,S}} Q_{0,S}, \qquad L^{\natural} \otimes
L^{\natural} \simeq L^{\flat} = L \wedge_{\{\pm 1\}} \mu_4^{*},
\]
$\mu_4^{*} = \mu_4 \setminus \mu_2$ étant vu comme revêtement étale quadratique
de $S_0$.\footnote{L'exposant que la transcription rend $\natural$ est, sur les
pages 79 à 96, un signe en zigzag en forme d'éclair, non identifié.} Pour un
revêtement étale cubique $\underline{J}$ de $S$, $\Sigma_{\underline{J}} =
\mathbb{P}(V_{\underline{J}}(\mathcal{O}_S))$, $V_{\underline{J}}(\mathcal{O}_S) =
\operatorname{Ker}(p_* \mathcal{O}_{\underline{J}} \xrightarrow{\operatorname{tr}}
\mathcal{O}_S)$, on pose $T_{0} = \mathcal{O}_{\Sigma_{\underline{J}}}(1)^{\natural}$,
de sorte que
\[
T_0^{\otimes 2} \simeq \mathcal{O}(2) \wedge_{\{\pm 1\}} \mu_4^{*} \simeq
T_{\Sigma_{\underline{J}}} \wedge_{\{\pm 1\}} (\mu_4^{*} \wedge
\underline{\omega}_{\underline{J}}).
\]
Pour une courbe elliptique $E/S$, $\underline{J} = ({}_2E)^{*}$, $E/\{\pm 1\}
\simeq \Sigma_{\underline{J}}$, $t$ l'image de l'origine, et $T_{E/S,0}^{\otimes 2}
\simeq T_{\Sigma_{\underline{J}}, t}$.

\emph{Définition.} Un épinglage de Legendre de $E/S$ est un isomorphisme $k :
T_{E/S,0} \simeq \mathcal{O}(1)^{\natural}_t$ dont le carré, lu à travers les
isomorphismes précédents, provient d'une section $\alpha$ de $\mu_4^{*} \wedge
\underline{\omega}_{\underline{J}}$, c'est-à-dire d'un isomorphisme $\alpha :
\mu_4^{*} \simeq \underline{\omega}_{\underline{J}}$ --- qui est déterminé par
$k$. Pour $\alpha$ fixé, les $k$ forment un torseur sous $\mu_2$ ; sans fixer
$\alpha$, ils forment un torseur $\mathcal{L}_{E/S}$ sous $\mu_4$, le
\emph{torseur des épinglages de Legendre}, avec
\[
\mathcal{L}_{E/S} / \mu_2 \simeq \mu^{*}_{4,S} \wedge
\underline{\omega}_{\underline{J}} .
\]
On peut aussi le décrire sans $\natural$ : les isomorphismes $k' : T_{E/S,0}
\simeq \mathcal{O}(1)_t$ dont le carré provient d'une section de
$\underline{\omega}_{\underline{J}}$ forment un $\mu_4$-torseur $Q_{E/S}$ --- un
schéma en carrés combinatoires, de schéma des diagonales
$\underline{\omega}_{\underline{J}}$ ---, et $\mathcal{L}_{E/S} \simeq Q_{E/S}
\wedge^{\mu_4} Q_{0,S}$. On aurait pu définir les épinglages de Legendre comme
les trivialisations de $Q_{E/S}$ ; on a préféré $\mathcal{L}_{E/S}$, mieux relié
aux épinglages de Jacobi d'échelon $4$.

\emph{Automorphismes.} Le champ des courbes munies d'un épinglage de Legendre
équivaut à celui des triplets $(\underline{J}, t, \alpha)$, et n'est pas rigide :
un automorphisme est un automorphisme de $\underline{J}$ compatible à $\alpha$,
donc une section de $\mathfrak{S}^{+}_{\underline{J}}$, et, s'il est non trivial
sur chaque fibre, une section $\rho$ de $\underline{\omega}_{\underline{J}}$
faisant de $\underline{J}$ un torseur sous $\mathbb{Z}/3\mathbb{Z}$ ; la
compatibilité à $t$ dit que $t$ est un point fixe de $\rho$. Ce sont les courbes
d'invariant $j = 0$, la courbe « hexagonale ».\footnote{« hexagonale » est une
lecture incertaine (un mot en \emph{h-}, terminé en \emph{-gonale}) ; c'est ce
que le sens appelle : les points fixes d'un $3$-cycle de $\{0, 1, \infty\}$ sont
les racines primitives sixièmes de l'unité, $\lambda = e^{\pm i\pi/3}$, $j = 0$,
réseau hexagonal.}

L'épinglage de \emph{Legendre--Jacobi d'échelon $2$}, couple $(k, u)$ avec $u :
\{0, 1, \infty\}_S \simeq \underline{J}$, tue ces automorphismes : le champ
correspondant est représentable par
\[
U_{0,3} \times_{S_0} \mu_4^{*}, \qquad U_{0,3} = \mathbb{P}^1_{S_0} \setminus
\{0, 1, \infty\},
\]
et la courbe universelle y a un fibré tangent à l'origine isomorphe à
$\mathcal{O}(1)^{\natural}$, qui devient $\mathcal{O}_{U_{0,3}}(1)$ puisque
$Q_{0}$ est trivial au-dessus de $\mu_4^{*}$.\footnote{La page note
$\Delta_{E/S}$, aux pages 81 à 86, le fibré tangent $T_{E/S,0}$ ; on écrit
$T$. Le second nom propre de la parenthèse de la page 85 n'est pas lu.}

\emph{Le groupe modulaire.} Les épinglages de Legendre forment un torseur sous
$\mu_4$, ceux de Jacobi d'échelon $2$ un torseur sous $\mathfrak{S}_3$ ; les
bi-épinglages, un torseur sous $\mu_4 \times \mathfrak{S}_3$, qui opère donc à
droite sur $U_{0,3} \times \mu_4^{*}$ : $\mathfrak{S}_3$ par l'action
tautologique sur $U_{0,3}$ et par la signature sur $\mu_4^{*}$, $\mu_4$
trivialement sur $U_{0,3}$ et par $\lambda \mapsto \lambda^2 \in \{\pm 1\}$ sur
$\mu_4^{*}$. Le sous-groupe qui opère trivialement sur $\mu_4^{*}$ est
\[
\mathbb{D}'_3 = \mathfrak{S}_3 \times_{\mu_2} \mu_4 = \{(g, \lambda) \mid
\mathrm{sg}(g) = \lambda^2\},
\]
extension $1 \to \{\pm 1\} \to \mathbb{D}'_3 \to \mathfrak{S}_3 \to 1$, d'ordre
$12$, d'abélianisé $\mu_4$ par le caractère $\chi_4$.\footnote{C'est le groupe
dicyclique d'ordre $12$, $\mathbb{Z}/3\mathbb{Z} \rtimes \mathbb{Z}/4\mathbb{Z}$
(le nom est le nôtre). La page note $\mathbb{D}_3 = \mathfrak{S}_3$ et
$\mathbb{D}'_3$ d'un D à double barre.} Au-dessus de $\widetilde{S}_0 =
\mu^{*}_{4,S_0}$, où une racine $i$ est choisie et $\mu_4 \simeq
\mathbb{Z}/4\mathbb{Z}$, on distingue les épinglages « directs », ceux dont la
section de $\mu_4^{*}$ est $i$ ; leur champ est représentable par
$U_{0,3,\widetilde{S}_0}$, sur lequel opère $\mathbb{D}'_3$.

Sur les fibrés, l'action modulaire $T_{g,\lambda}$ tord l'action tautologique :
par $\lambda$ sur $\mathcal{O}_{U_{0,3}}(1)$, par $\lambda^2$ sur
$\mathcal{O}(2)$, par $\lambda^2 \mathrm{sg}(g)$ sur le fibré tangent. Sur
$\mathbb{D}'_3$, où $\lambda^2 = \mathrm{sg}(g)$, il reste
\[
T_g \zeta = \chi_4(g)\,(g \cdot \zeta), \qquad T_g \eta = \mathrm{sg}(g)\,(g
\cdot \eta), \qquad T_g \tau = g \cdot \tau,
\]
pour $\zeta$, $\eta$, $\tau$ sections respectives de $\mathcal{O}(1)$,
$\mathcal{O}(2)$ et du fibré tangent : le fibré tangent porte son action
naturelle.\footnote{Dans la formule (52 bis), l'exposant $^{0}$ de
$\mathrm{sg}(g^{0})$ est lu tel qu'écrit ; on entend la signature de l'image de
$g$ dans $\mathfrak{S}_3$.}

Reste, écrit-il ensuite, à définir les épinglages de Jacobi d'échelon $4$
relatifs à ceux d'échelon $2$ --- sections de ${}_4E \to {}_2E$ au-dessus de
$\underline{J}_E$, avec $\xi_i + \xi_j + \xi_k = 0$ sur chaque fibre
géométrique, ce que la partie II, §4 appelait épinglage relatif --- et à
\emph{montrer comment un tel épinglage définit canoniquement un épinglage de
Legendre}. C'est le programme de la fin de la partie IV ; il n'est pas mené à
terme dans le dossier.

\pagerange{93}{95}
\section*{La courbe standard $E_\alpha$ (pages 93 à 95)}

Pour $\underline{J}$ et $t$ donnés, les courbes correspondantes équivalent aux
couples $(T, T^{\otimes 2} \simeq T_{\Sigma,t})$, et, un tel $T_0$ étant fixé,
aux $\mu_2$-torseurs $\underline{\varepsilon}$, par $T = T_0
\wedge^{\mu_2} \underline{\varepsilon}$. Avec $T_0 = \mathcal{O}_\Sigma(1)(t)^{\natural}$,
pour avoir un isomorphisme « raisonnable » $T_0^{\otimes 2} \simeq T_{\Sigma,t}$
il faut un $\alpha : \mu_4^{*} \simeq \underline{\omega}_{\underline{J}}$. Il
définit une courbe $E_\alpha$ munie d'un épinglage de Legendre standard
subordonné à $\alpha$ --- la courbe \emph{universelle} pour ces données ---, et
pour toute courbe $E$ de mêmes $\underline{J}$, $t$, on a un isomorphisme de
$\mu_2$-torseurs
\[
\underline{\operatorname{Isom}}_{S, \underline{J}}(E_\alpha, E)
\xrightarrow{\;\sim\;} \mathcal{L}_\alpha(E), \qquad u \longmapsto u(k_\alpha),
\]
l'indice $\underline{J}$ signifiant qu'on prend les isomorphismes compatibles à
$\underline{J} \simeq {}_2E_\alpha \simeq {}_2E$. Avec une section $\omega$ de
$\underline{\omega}_{\underline{J}}$ au lieu de $\alpha$, on définirait de même
une courbe $E_\omega$ munie d'un épinglage de Legendre « naïf ». Les deux points
de vue coïncident au-dessus de $\widetilde{S}_0 = \mu^{*}_{4,S_0}$.

\pagerange{95}{99}
\section*{Explicitation calculatoire (pages 95 à 99)}

Sur un corps algébriquement clos $k$, on décrit $E$ par $J$, le plan $V_J =
\operatorname{Ker}(k^J \to k)$, $\Sigma_J = \mathbb{P}(V_J)$ avec ses trois
points $\xi_i = k\,(e_j - e_k)$, une droite $F_t \subset V_J$ (le point $t$), la
droite $\underline{t}_E = T_{E,0}$ et un isomorphisme
\[
\nu : \underline{t}_E^{\otimes 2} \simeq T_{\Sigma,t} \simeq
\mathcal{O}_t(2)(\underline{\omega}) \simeq F_t^{\otimes(-2)}(\underline{\omega}),
\qquad \text{ou} \qquad \check\nu : \underline{t}_E^{\otimes(-2)} \simeq
F_t^{\otimes 2}(\underline{\omega}),
\]
puisque $\mathcal{O}_t(1) = V_J / F_t \simeq \check F_t(\underline{\omega})$.
Un épinglage de Legendre $k : \underline{t} \simeq \mathcal{O}_t(1)^{\natural}$
est une famille $(k_\iota)_{\iota \in \mu_4^{*}}$ d'isomorphismes $k_\iota :
\underline{t} \simeq \check F_t(\underline{\omega})$ avec $k_{\iota'} = \iota\,
k_\iota$, ou, dualement, $\check k_\iota : \check{\underline{t}} \simeq
F_t(\underline{\omega})$ avec $\check k_{\iota'} = \iota^{-1} \check k_\iota$,
telle que $k_\iota^{\otimes 2}(\delta^{\otimes 2}) = \nu(\delta^{\otimes 2})
\wedge \omega_\iota$ pour un $\omega_\iota \in \underline{\omega}$ déterminé par
$k_\iota$. On a $\omega_{\iota'} = $ l'autre élément de $\underline{\omega}$,
donc $\iota \mapsto \omega_\iota$ est la bijection $\alpha = \alpha_k : \mu_4^{*}
\simeq \underline{\omega}$.

Écrivons $\check k_\iota(\check\delta) = \psi^{\omega_\iota}(\delta) \wedge
\omega_\iota$, $\psi^{\omega}(\delta) \in F_t$, pour $\delta \in
\underline{t}^{*}$. Les conditions deviennent
\[
\psi^{\omega'}(\delta) = \iota_\omega\, \psi^{\omega}(\delta), \quad \iota_\omega
= \alpha^{-1}(\omega) ; \qquad \psi^{\omega}(\lambda\delta) = \lambda^{-1}
\psi^{\omega}(\delta) ; \qquad \psi^{\omega}(\delta)^{\otimes 2} =
\check\nu(\check\delta^{\otimes 2}) \wedge \omega,
\]
la première venant de $-1/\iota = \iota$ dans $\mu_4^{*}$.

\begin{quote}
\textbf{Lemme} (page 98). Les épinglages de Legendre de données
$(J, F_t \subset V_J, \underline{t}, \check\nu)$ s'identifient aux systèmes
$\bigl(\alpha, (\psi^{\omega}(\delta))_{\omega \in \underline{\omega}, \delta
\in \underline{t}^{*}}\bigr)$, $\alpha : \mu_4^{*} \simeq \underline{\omega}$,
$\psi^{\omega}(\delta) \in F_t$, satisfaisant ces trois conditions.
\end{quote}

En coordonnées, $\psi^{\omega}(\delta) = \sum_i \psi^{\omega}_i(\delta)\, e_i$
avec $\sum_i \psi^{\omega}_i(\delta) = 0$,
$\psi_i^{\omega'}(\delta) = \iota_\omega \psi_i^{\omega}(\delta)$,
$\psi_i^{\omega}(\lambda\delta) = \lambda^{-1} \psi_i^{\omega}(\delta)$ ; si
$F_t = \{\sum \lambda_i e_i \mid \sum \lambda_i = 0, \sum c_i \lambda_i = 0\}$,
avec $(c_i)$ déterminé à un scalaire et à un vecteur diagonal près,
l'appartenance à $F_t$ s'écrit $\sum_i c_i \psi^{\omega}_i(\delta) = 0$, et il
suffit de l'écrire pour un seul couple $(\omega, \delta)$. « Il reste ensuite à
exprimer » la dernière condition : la page s'arrête là.\footnote{La page 98
porte une note marginale : moyennant la première et la deuxième condition, les
deux relations quadratiques, pour les deux $\omega$, sont équivalentes. Le
numéro (77) est employé deux fois sur la page 99, dont une fois dans un passage
biffé.}

\pagerange{100}{108}
\section*{IV. Descriptions explicites par fonctions $\wp$ de Weierstrass --- L'algèbre des fonctions à pôle à l'origine (pages 100 à 108)}

Soit $E$ une courbe elliptique sur un corps algébriquement clos $k$ de
caractéristique $\neq 2$. Posons $\mathcal{F}_n = \Gamma(E, \mathcal{O}(n \cdot
0_E))$ : par Riemann--Roch, $\dim \mathcal{F}_n = n$ pour $n \geq 1$, et
$\mathcal{F}_0 = \mathcal{F}_1 = k$. Les $\mathcal{F}_n$ forment une algèbre
graduée $\mathcal{F}_{*}$ dont le Proj est $E$, et $\mathcal{F}_\infty =
\varinjlim \mathcal{F}_n = \Gamma(E \setminus 0_E, \mathcal{O})$. Pour $n \geq
2$, le coefficient dominant donne une suite exacte $0 \to \mathcal{F}_{n-1} \to
\mathcal{F}_n \xrightarrow{\varepsilon_n} \underline{t}_E^{\otimes n} \to 0$,
$\underline{t}_E = T_{E, 0_E}$ : si $z$ est un paramètre local dual de $\delta$,
$f = c\, z^{-n} + \cdots$ a pour image $c\, \delta^{\otimes n}$.

Pour $\delta \in \underline{t}_E$, la dérivation invariante $\Theta_\delta$ est
une dérivation de degré $+1$ de $\mathcal{F}_{*}$, et $\varepsilon_{n+1}(\Theta_\delta
f) = -n\, \delta \otimes \varepsilon_n(f)$.\footnote{La page fait de la flèche
de droite de son diagramme (9) la multiplication par $\delta$ ; la dérivée de
$z^{-n}$ est $-n z^{-n-1}$, d'où le facteur $-n$. Il est sans conséquence pour
$n = 2$ (caractéristique $\neq 2$), mais non pour la suite.} Pour $n$ inversible
dans $k$, cette flèche est injective, d'où $\operatorname{Ker} \Theta_\delta
\cap \mathcal{F}_n = k$ : en caractéristique nulle, $\operatorname{Ker}
\Theta_\delta = k$.\footnote{La page affirme $\operatorname{Ker} \Theta_\delta =
\mathcal{F}_1 = k$ sans restriction. En caractéristique $p > 0$, une dérivation
tue les puissances $p$-ièmes, et $\operatorname{Ker} \Theta_\delta$ contient
$\mathcal{F}_{*}^{\,p}$. Le résultat n'est pas utilisé dans la suite.}
L'involution $f \mapsto \check f$, $\check f(x) = f(-x)$, anticommute à
$\Theta_\delta$ et scinde $\mathcal{F}_n = \mathcal{F}_n^{+} \oplus
\mathcal{F}_n^{-}$, avec les règles de parité évidentes. On a
$\mathcal{F}^{+}_{2n} \simeq \mathcal{F}^{+}_{2n+1} = \Gamma(\Sigma_J, n\{t\})$,
$\Sigma_J = E/\{\pm 1\}$ ; $\mathcal{F}_n^{-} = 0$ pour $n \leq 2$ ;
$\mathcal{F}_3^{-} \simeq \underline{t}_E^{\otimes 3}$ ; et
$\mathcal{F}^{-}_{2n+1} = \mathcal{F}^{-}_{2n+2} = \mathcal{F}_3^{-} \cdot
\mathcal{F}^{+}_{2n-2}$.

Choisissons $\delta \neq 0$ et $\wp \in \mathcal{F}_2$ avec $\varepsilon_2(\wp)
= \delta^{\otimes 2}$, et posons
\[
y = -\tfrac12\, \Theta_\delta \wp, \qquad \text{de sorte que } \varepsilon_3(y)
= \delta^{\otimes 3} .
\]
\footnote{La page pose $\wp' = \Theta_\delta(\wp)$ et écrit
$\varepsilon_3(\wp') = \delta^{\otimes 3}$ ; en fait $\varepsilon_3(\Theta_\delta
\wp) = -2\,\delta^{\otimes 3}$, et avec son $\wp'$ la relation (28) s'écrirait
$\wp'^2 = 4(\wp^3 + a_1 \wp^2 + a_2 \wp + a_3)$ --- c'est la relation classique
$\wp'^2 = 4\wp^3 - g_2 \wp - g_3$. Toute la suite de la page (propositions 2 et
3, formules d'addition, points d'ordre $4$) n'utilise que la relation unitaire ;
elle vaut telle quelle pour $y = -\frac12 \Theta_\delta \wp$, qu'on écrit à la
place de son $\wp'$.} Alors $1, \wp$ est une base de $\mathcal{F}_2 =
\mathcal{F}_3^{+}$, $y$ une base de $\mathcal{F}_3^{-}$, et plus généralement
$1, \wp, \ldots, \wp^n$ et $y, y\wp, \ldots, y\wp^{n-1}$ sont des bases de
$\mathcal{F}^{+}_{2n}$ et de $\mathcal{F}^{-}_{2n+1}$.

\emph{Proposition 1.} Le noyau de la restriction $\mathcal{F}_3 \to k^J$ aux
trois points d'ordre $2$ est $\mathcal{F}_3^{-} = k\, y$ : une fonction impaire
s'annule aux points d'ordre $2$, et la restriction est injective sur
$\mathcal{F}_2$. \emph{Corollaire} : la matrice de Vandermonde
$(\wp(\varepsilon_i)^\alpha)_{i \in J, 0 \leq \alpha \leq 2}$ est inversible,
c'est-à-dire que les trois valeurs $\wp(\varepsilon_i)$ sont distinctes ; on le
voit en normalisant $J = \{0, 1, \infty\}$, $\wp = 1/(z - t)$, $t \in
U_{0,3}(k)$.\footnote{La page écrit $t \in k \setminus \{0, 1\} = U_{0,1}(k)$ ;
c'est $\mathbb{P}^1(k) \setminus \{0, 1, \infty\} = U_{0,3}(k)$.}

Comme $\varepsilon_6(y^2) = \varepsilon_6(\wp^3) = \delta^{\otimes 6}$ et que $y^2
- \wp^3$ est pair, il vient
\[
y^2 = \wp^3 + a_1 \wp^2 + a_2 \wp + a_3, \qquad a_1, a_2, a_3 \in k,
\]
uniquement déterminés pour $\delta$ choisi. \emph{Proposition 2} : les $a_\alpha$
sont caractérisés par l'annulation de $\wp^3 + a_1 \wp^2 + a_2 \wp + a_3$ sur
$J$, puisque $y|_J = 0$ et que $1, \wp, \wp^2$ sont indépendants sur $J$.
\emph{Corollaire 1} : $\wp^3 + a_1\wp^2 + a_2\wp + a_3 = \prod_{i \in J} (\wp -
\wp(\varepsilon_i))$. \emph{Corollaire 2} : $a_1 = 0$ si et seulement si $\sum_i
\wp(\varepsilon_i) = 0$ ; si la caractéristique est $\neq 3$, ces $\wp$ forment
un supplémentaire $W_E \simeq \underline{t}_E^{\otimes 2}$ de $\mathcal{F}_1$
dans $\mathcal{F}_2$, dont les éléments non nuls sont les fonctions de
Weierstrass. \emph{Corollaire 3} :
\[
\mathcal{F}_\infty \simeq k[P, Y] / \bigl(Y^2 - (P^3 + a_1 P^2 + a_2 P +
a_3)\bigr), \qquad P \mapsto \wp, \ Y \mapsto y,
\]
l'involution étant $Y \mapsto -Y$, et $\Theta_\delta$ étant induite par la
dérivation $\Theta_\delta P = -2Y$, $\Theta_\delta Y = -(3P^2 + 2a_1 P + a_2)$,
qui respecte l'idéal --- on l'obtient en dérivant l'équation.\footnote{Avec son
$\wp' = \Theta_\delta \wp$ et sa relation unitaire, la page obtient $\wp'' =
\frac12(3\wp^2 + 2a_1\wp + a_2)$ ; dans la normalisation correcte,
$(\Theta_\delta)^2 \wp = 2(3\wp^2 + 2a_1\wp + a_2)$.} \emph{Remarque} : dans les
bases ci-dessus, on explicite la multiplication et la dérivation en termes de
$a_1, a_2, a_3$. Inversement, si le discriminant de $Z^3 + a_1Z^2 + a_2Z + a_3$
est non nul, ces formules définissent une courbe elliptique, revêtement double
$\wp : E \to \mathbb{P}^1_k \simeq E/\{\pm 1\}$ ramifié à l'origine (au-dessus
de $\infty$) et en ses trois racines, et l'inclusion $J \subset \mathbb{P}^1_k$
identifie $\Sigma_J$ à $\mathbb{P}^1_k$, $t \mapsto \infty$.

\pagerange{109}{112}
\section*{La loi d'addition (pages 109 à 112)}

Un point $x \neq 0$ est la donnée de $x_0 = \wp(x)$ et d'une racine $x' = y(x)$
de $X'^2 = x_0^3 + a_1x_0^2 + a_2x_0 + a_3$ (deux racines distinctes si $x_0
\notin D_J$).

\begin{quote}
\textbf{Proposition 3.} Pour $x, y, z \in E(k) \setminus \{0\}$ deux à deux
distincts, $x + y + z = 0$ si et seulement si
\[
\det \begin{pmatrix} 1 & \wp(x) & y(x) \\ 1 & \wp(y) & y(y) \\ 1 & \wp(z) & y(z)
\end{pmatrix} = 0 .
\]
\end{quote}
En effet $x + y + z = 0$ équivaut à $\{x\} + \{y\} + \{z\} \sim 3\{0\}$, donc à
l'existence d'une $f = y + \alpha \wp + \beta \in \mathcal{F}_3$ de diviseur
$\{x\} + \{y\} + \{z\} - 3\{0\}$ ; cela donne la dépendance des trois colonnes.
Inversement, trois points distincts n'ont pas la même image par $\wp$ (les
fibres ont deux points), donc les colonnes $(1)$ et $(\wp)$ sont indépendantes
et la relation fournit un tel $f$, qui s'annule aux trois points distincts et
n'a pas d'autre zéro.\footnote{Les trois notes marginales de la page 109,
écrites de biais, ne se lisent que par fragments ; l'une remarque que la
condition force $\wp(x), \wp(y), \wp(z)$ deux à deux distincts.}

Pour calculer $z = -(x + y)$, avec $x_0 \neq y_0$ : la droite $y = \lambda \wp +
\mu$, $\lambda = \frac{x' - y'}{x_0 - y_0}$, $\mu = \frac{y'x_0 - x'y_0}{x_0 -
y_0}$, coupe la cubique en trois points dont $x_0$ et $y_0$, d'où par la somme
des racines
\[
z_0 = -x_0 - y_0 + \Bigl( \frac{x' - y'}{x_0 - y_0} \Bigr)^2 - a_1, \qquad z' =
\lambda z_0 + \mu .
\]
Établie d'abord sous les conditions $x \neq y$, $y \neq -2x$, $x \neq -2y$, la
formule vaut par prolongement dès que $\wp(x) \neq \wp(y)$, c'est-à-dire $y
\notin \{\pm x\}$. C'est la loi d'addition par la corde.

\pagerange{113}{117}
\section*{Les involutions $\sigma_i$ et les points d'ordre $4$ (pages 113 à 117)}

$\wp$ induit $E/\{\pm 1\} \simeq \mathbb{P}^1_k$ avec $\wp(0_E) = \infty$, ce qui
la détermine à une transformation affine près ; la normalisation de Weierstrass
$\sum \zeta_i = 0$, $\zeta_i = \wp(\varepsilon_i)$, possible en caractéristique
$\neq 3$, la détermine à un facteur près. En général $a_1 = -\sum \zeta_i$, $a_2
= \sum \zeta_i \zeta_j$, $a_3 = -\prod \zeta_i$. L'involution $\sigma_i$ de
$\mathbb{P}^1_k$ induite par la translation par $\varepsilon_i$ échange $\zeta_i
\leftrightarrow \infty$ et $\zeta_j \leftrightarrow \zeta_k$ :
\[
\sigma_i Z = \frac{\zeta_i Z + b_i}{Z - \zeta_i}, \qquad b_i = \zeta_j\zeta_k -
\zeta_i(\zeta_j + \zeta_k) = 2\zeta_j\zeta_k - a_2 .
\]
Ses points fixes sont les racines $Z_i$ de $Z^2 - 2\zeta_i Z - b_i = 0$, et
\[
(Z_i - \zeta_i)^2 = \zeta_i^2 - a_2 + 2\zeta_j\zeta_k = (\zeta_i - \zeta_j)(\zeta_i
- \zeta_k) .
\]
\footnote{La forme factorisée est la nôtre ; c'est la formule classique de la
division par $2$. La page passe par $b_i = -(a_2 + 2a_3/\zeta_i)$, qui suppose
$\zeta_i \neq 0$ --- ce qui l'inquiète dans deux notes marginales --- et arrive à
$(Z_i - \zeta_i)^2 = -\frac{1}{\zeta_i}(a_1\zeta_i^2 + 2a_2\zeta_i + 3a_3)$,
juste. Mais en recopiant cette expression page 117 elle la réécrit $-(a_1\zeta_i
+ 2a_2 + 3\zeta_j\zeta_k)$, où il faut $-3\zeta_j\zeta_k$, puisque $a_3/\zeta_i =
-\zeta_j\zeta_k$ ; et la note marginale de la page 117, $Z_i^2 = 2\zeta_i Z_i -
(a_2 + 2\zeta_j\zeta_k)$, doit se lire $Z_i^2 = 2\zeta_i Z_i - (a_2 -
2\zeta_j\zeta_k)$. Le renvoi « (66) » de la page 115 se lit (60).} Les douze
points d'ordre $4$ sont au-dessus des six points $Z_i$ : choisir $x_i$ avec $2x_i
= \varepsilon_i$ revient à choisir $\wp(x_i) = Z_i$ (deux choix), puis $y(x_i) =
Z'_i$ avec
\[
Z_i'^2 = (Z_i - \zeta_i)(Z_i - \zeta_j)(Z_i - \zeta_k)
\]
(deux choix) : quatre choix, comme il se doit. Les points $x_i$ étant deux à
deux distincts et non nuls, la proposition 3 traduit la condition de Jacobi
relative d'échelon $4$, $\sum_{i \in J} x_i = 0$, en
\[
\det \begin{pmatrix} 1 & Z_1 & Z'_1 \\ 1 & Z_2 & Z'_2 \\ 1 & Z_3 & Z'_3
\end{pmatrix} = 0, \qquad \text{i.e.} \qquad \varphi_1 + \varphi_2 + \varphi_3 =
0, \quad \varphi_i = Z_j Z'_k - Z_k Z'_j,
\]
où $(i, j, k)$ parcourt les permutations circulaires, ou, pour un ordre
circulaire $\omega$, $j = \omega i$, $k = \omega^2 i$.

\pagerange{117}{125}
\section*{Vers l'épinglage de Legendre (pages 117 à 125)}

Grothendieck voudrait associer canoniquement à une telle donnée $(x_i)_{i \in
J}$ un épinglage de Legendre de $E$, fonctoriel en $E$ et compatible aux
changements de base. Fixons $\iota \in \mu_4^{*}(k)$ ; la donnée $\alpha$
devient un ordre circulaire $\omega$ sur $J$ ($\alpha = \iota \wedge \omega$), et
la donnée $k$ un isomorphisme $K : \underline{t}_E \simeq \mathcal{O}_t(1)$ dont
le carré, lu à travers $\underline{t}_E^{\otimes 2} \simeq T_{\Sigma_J,t}$ et
$\mathcal{O}_t(2) \simeq \mathcal{O}_t(2)(\underline{\omega}) \simeq
T_{\Sigma_J,t}$ (via $\alpha$), soit l'identité : $K$ est déterminé au signe près
pour $\alpha$ fixé, à $\mu_4$ près sinon.

La suite exacte $0 \to F_t \to V_J \to \mathcal{O}_t(1) \to 0$, avec $F_t \simeq
\mathcal{O}_t(-1)(\underline{\omega})$ et $V_J \simeq \Gamma(\Sigma_J,
\mathcal{O}(1))$, tensorisée par $\check F_t \simeq
\mathcal{O}_t(1)(\underline{\omega})$, devient canoniquement
\[
0 \to k \cdot 1 \to \Gamma(\Sigma_J, \mathcal{O}(\{t\})) \to T_{\Sigma_J,t} \to
0,
\]
où le terme médian, les fonctions sur $\Sigma_J$ à pôle simple en $t$, est
$\mathcal{F}_2$. Se donner $K$ équivaut à se donner $\check K :
\check{\underline{t}}_E \simeq \mathcal{O}_t(-1) \simeq
\mathcal{O}_t(-1)(\underline{\omega}) \simeq F_t$ (via $\omega$), donc un
monomorphisme $\check{\underline{t}}_E \hookrightarrow V_J \subset k^J$ d'image
$F_t$, donc une famille $(\delta_i)_{i \in J} \in \underline{t}_E^J$ avec a)
$\sum \delta_i = 0$, b) une relation de proportionnalité exprimant que l'image
est $F_t$, c) la relation de l'épinglage de Legendre. Écrivant $\delta_i =
\varphi_i^{\omega}(\delta)\, \delta$ pour $\delta \in \underline{t}_E^{*}$ :
\[
\varphi_i^{\omega}(\lambda\delta) = \lambda^{-1} \varphi_i^{\omega}(\delta),
\qquad \sum_i \varphi_i^{\omega}(\delta) = 0, \qquad k \cdot
(\varphi_i^{\omega}(\delta))_i = F_t, \qquad \varphi_i^{\omega'}(\delta) =
\iota_\omega \varphi_i^{\omega}(\delta),
\]
et la condition b) les détermine à un scalaire près, c) au signe près. « Et la
condition c) s'écrit\ldots\ Dieu sait comment --- il faudra y revenir » ; la
page 120 s'arrête sur « Posons ».\footnote{Les numéros (75) et (83) sont portés
deux fois aux pages 117 et 119 ; la transcription les garde, et on n'en cite
aucun.}

\emph{La droite $F_t$} (pages 122 à 125). Trois feuillets pris en largeur, sans
titre, font le calcul qui donne la condition b) explicitement. Pour $\xi_1, \xi_2,
\xi_3 \in k$ distincts, le système $\sum \lambda_i = 0$, $\sum \lambda_i \xi_i =
0$ a pour solution $\lambda = (\xi_2 - \xi_3, \xi_3 - \xi_1, \xi_1 - \xi_2)$, et la
base $(f_1, f_2)$ de $V_J$ définie par $\xi_i f_1 + f_2 = (e_j - e_k)/(\xi_j -
\xi_k)$ est
\[
\Delta_0 f_1 = \sum_{i} (\xi_j - \xi_k)\, e_i, \qquad -\Delta_0 f_2 = \sum_i
\xi_i (\xi_j - \xi_k)\, e_i, \qquad \Delta_0 = (\xi_1 - \xi_2)(\xi_2 -
\xi_3)(\xi_3 - \xi_1),
\]
sommes sur les permutations circulaires $(i, j, k)$. L'isomorphisme $k^2 \simeq
V_J$, $(u, v) \mapsto u f_1 + v f_2$, envoie le point $\xi_i$ de la droite sur
la droite $k\,(e_j - e_k)$ de $V_J$, c'est-à-dire sur le point $i$ de
$\Sigma_J$, et le point $\infty$ sur $k f_1$. Avec $\xi_i = \wp(\varepsilon_i)$,
le point $t = \wp(0_E) = \infty$ correspond donc à
\[
F_t = k \cdot \sum_i \bigl(\wp(\varepsilon_j) - \wp(\varepsilon_k)\bigr)\, e_i ,
\]
de sorte que la condition b) s'écrit $\varphi_i^{\omega}(\delta) = c\,(\zeta_j -
\zeta_k)$, $(i, j, k)$ circulaire.\footnote{Le rattachement des pages 122 à 125
à la condition b) de la page 120 est le nôtre ; il est appelé par la dernière
ligne de la page 125, « $t$ correspond donc : $\sum (\xi_j - \xi_k)\, e_i =
\sum_i (\wp(\varepsilon_j) - \wp(\varepsilon_k))\, e_i$ ». Les pages 121
(blanche) et 123 (un avis de soutenance dactylographié) s'intercalent sans
rapport.}

\pagerange{126}{140}
\section*{V. Approche transcendante --- première mouture (pages 126 à 140)}

Deux feuillets de titre précèdent cette partie : « Relations entre $M_{11}$ et
$M_{04}$ » (page 127), sous lequel le dossier ne contient rien, et « Approche
transcendante (21 p) » (page 128).\footnote{Les pages 126 (une annonce de
séminaire dactylographiée) et 129 (blanche) ne portent pas de mathématiques. Les
« 21 p » annoncées correspondent aux pages 130 à 151.} La page 130 est titrée
« Modules des courbes elliptiques : théorie transcendante ».

Une courbe elliptique sur $\mathbb{C}$ est la donnée de la droite complexe $V =
\underline{t}_E$ et d'un réseau $\Pi \subset V$ ; $\Pi$ redonne $V = \Pi \otimes
\mathbb{R}$ comme espace réel et $E = V/\Pi$ comme groupe de Lie réel, la
structure complexe venant de celle de $V$. Une structure complexe sur un espace
réel $V$ équivaut à un sous-espace complexe $L \subset V_{\mathbb{C}}$ tel que
$\Re : L \to V$ soit un isomorphisme réel ; on munit $V$ de la structure
\emph{conjuguée} de celle transportée de $L$, $u x = \varphi(-i\, \varphi^{-1}(x))$
--- choix justifié par la réciproque : pour $V$ complexe, le noyau $L'$ de
$V_{\mathbb{C}} \to V$, $x + iy \mapsto x + uy$, est $\{x + iux\}$, et $\Re :
L' \to V$ est antilinéaire.\footnote{La page écrit $u x = \varphi(-i\,
\varphi^{-1}(x)) = -\varphi(i\,\varphi(x))$ ; la seconde expression est un lapsus
pour $-\varphi(i\, \varphi^{-1}(x))$.} Variante : une structure complexe équivaut
à un sous-groupe compact $\mathbb{U} \subset \operatorname{Aut}_{\mathbb{R}}(V)$
isomorphe au cercle, avec ses deux sous-espaces propres dans $V_{\mathbb{C}}$,
plus le choix d'un de ses deux éléments d'ordre $4$, c'est-à-dire d'une
orientation.

Pour $V$ de rang réel $2$, les droites complexes de $V_{\mathbb{C}}$ sont les
points de $\Sigma = \mathbb{P}(V)(\mathbb{C})$, et $L$ définit une structure
complexe si et seulement si $L \cap V = 0$, c'est-à-dire si le point n'est pas
réel. Les structures complexes sur $V$ sont donc les points de $\Sigma \setminus
\Sigma_{\mathbb{R}}$, qui a deux composantes $\Sigma^{*}_\omega$, une par
orientation de $V$. Une courbe elliptique est ainsi la donnée a) d'un
$\mathbb{Z}$-module libre $\Pi$ de rang $2$ ($= \pi_1(E, 0)$) ; b) d'une
orientation $\omega$ de $\Pi_{\mathbb{R}}$, c'est-à-dire d'une base de
$\bigwedge^2 \Pi$ ; c) d'un point de $\Sigma^{*}_\omega(V)$.

Une \emph{rigidification de Torelli} est un isomorphisme $u : \mathbb{Z}^2
\simeq \Pi$ tel que $z_1 \wedge z_2 = \omega$, $z_i = u(e_i)$. Il identifie
$\Sigma^{*}_{V,\omega}$ au demi-plan $\mathfrak{D} = \{z \mid \Im z > 0\}$, le
point $z$ correspondant à la droite $L_z = \mathbb{C} \cdot (z e_1 + e_2)$. Les
classes d'isomorphie de courbes avec rigidification de Torelli sont en
bijection avec $\mathfrak{D}$, et $\operatorname{SL}(2, \mathbb{Z})$ opère à
droite par $u \mapsto u \circ g$. Le calcul : $\Re(\lambda(z e_1 + e_2))$ donne
$\lambda_1 = -i/y$, $\lambda_2 = 1 + ix/y$ comme antécédents de $e_1, e_2$, de
quotient $-x + iy$ ; comme la structure de $V$ est la conjuguée, $z_2/z_1 =
-x - iy = -z$, soit
\[
z = -(z_2 / z_1) .
\]
Alors $u \circ g$ correspond à $z' = \frac{dz - b}{-cz + a}$, et l'action à gauche
$u \mapsto u \circ g^{-1}$ à
\[
z' = \frac{az + b}{cz + d}
\]
--- « \emph{victoire}, c'est la formule habituelle ». Mais on prendra garde que le
nombre attaché à la base $(z_1, z_2)$ n'est pas $z_2/z_1$, comme il paraîtrait
naturel : repérer par $\tau = z_2/z_1$ donnerait une autre formule de
transformation.\footnote{La page 138 écrit $z = -(z_2/z_1)^{-} =
-\bar z_2/\bar z_1$, ce qui contredit $z = -(z_2/z_1)$ de la page 136 par une
conjugaison, sauf à calculer le quotient dans $L$ et non dans $V$ : c'est le
quotient $\lambda_2/\lambda_1 = -\bar z$. Une longue note marginale de la page
138, en partie biffée, reprend la question des conventions, et envisage
$\mathbb{C} \cdot (1, z)$ au lieu de $\mathbb{C} \cdot (z, 1)$.}

\begin{quote}
\textbf{Théorème} (page 138). Le demi-plan $\mathfrak{D}$, muni de la famille
analytique de courbes elliptiques $E_z = \mathbb{C}/(\mathbb{Z} + \mathbb{Z}z)$,
est un revêtement universel du champ analytique $M_{1,1}^{\mathrm{an}}$,
correspondant à la classification des courbes à rigidification de Torelli ; son
groupe d'automorphismes au-dessus de $M_{1,1}^{\mathrm{an}}$ est
$\operatorname{SL}(2, \mathbb{Z})$, opérant par $gz = \frac{az + b}{cz + d}$, ce
qui correspond à l'action à gauche $g \cdot u = u \circ g^{-1}$ sur les
structures de Torelli.
\footnote{La page écrit $E_z = \mathbb{C}/(\mathbb{Z} \cdot 1 + \mathbb{Z} \cdot
(-\bar z))$. Cette famille dépend anti-holomorphiquement de $z$ ; avec $z_2/z_1 =
-z$ dans $\underline{t}_E$, le réseau est $\mathbb{Z} + \mathbb{Z}(-z) =
\mathbb{Z} + \mathbb{Z}z$. « Revêtement universel » s'entend du champ (ou de
l'orbifold) : $-1 \in \operatorname{SL}(2, \mathbb{Z})$ agit trivialement sur
$\mathfrak{D}$ mais par $-1$ sur les courbes, et l'espace des modules grossier,
la droite des $j$, est simplement connexe.}
\end{quote}

\emph{Corollaire 1.} $\mathcal{T}_{1,1} = \pi_1(M_{1,1}) \simeq
\operatorname{SL}(2, \mathbb{Z})$, calculé relativement au revêtement universel
$\mathfrak{D}$, et $M_{1,1} \simeq [\mathfrak{D} / \operatorname{SL}(2,
\mathbb{Z})]$. « Utilisant les grands théorèmes orthodoxes » --- la théorie de
Teichmüller, l'espace de Teichmüller étant ici $\mathfrak{D}$ :
\emph{Corollaire 2} : pour une surface compacte orientée de genre $1$ pointée
$(E, a)$, le groupe $\mathcal{T}$ des classes d'isotopie d'automorphismes
directs de $(E, a)$ s'envoie isomorphiquement sur
$\operatorname{Aut}^{+}_{\mathbb{Z}}(\Pi)$, $\Pi = \pi_1(E, a)_{\mathrm{ab}} =
H_1(E, \mathbb{Z})$. \emph{Corollaire 3} : pour un groupe $\pi$ de type
$(1, 1)$, $H = \pi_{\mathrm{ab}}$, libre de rang $2$, le groupe
$\operatorname{Autext}^{+}(\pi)$ des automorphismes extérieurs compatibles à
l'orientation s'envoie isomorphiquement sur $\operatorname{Aut}^{+}_{\mathbb{Z}}(H)
\simeq \operatorname{SL}(2, \mathbb{Z})$.\footnote{La page écrit « de rang 1 »
pour $H$ : l'abélianisé du groupe fondamental d'un tore épointé, groupe libre à
deux générateurs, est de rang $2$. Le corollaire 3 est un énoncé de Nielsen
(pour $\operatorname{Out}(F_2) \simeq \operatorname{GL}(2, \mathbb{Z})$) ; la
page ne cite personne. La notation $\mathcal{T}_{1,1}$ est celle des groupes de
Teichmüller $\mathcal{T}_{g,\nu}$.}

La page 140 commence le calcul du fibré tangent le long de la section unité de
la courbe universelle rigidifiée : comme fibré réel, c'est le fibré sur
$\mathfrak{D}$ défini par l'action standard de $\operatorname{SL}(2, \mathbb{Z})$
sur $\mathbb{R}^2$ ; le fibré principal associé s'obtient en remplaçant
$\mathbb{R}^2$ par $(\mathbb{R}^2)^{*} \simeq \mathbb{C}^{*}$, dont un revêtement
universel est $z \mapsto \exp(2i\pi z)$. La page s'arrête là.

\pagerange{141}{151}
\section*{Approche transcendante --- seconde mouture (pages 141 à 151)}

La seconde rédaction évite la structure conjuguée en prenant directement $\tau =
z_2/z_1$. Une courbe elliptique sur $\mathbb{C}$ est déterminée par le réseau $H
= \pi_1(E, 0) \simeq H_1(E, \mathbb{Z})$ et une structure complexe sur $H
\otimes \mathbb{R} \simeq \underline{t}_E$ ; sur une base analytique $S$, par un
système local $\mathcal{H}$ de $\mathbb{Z}$-modules libres de rang $2$ et une
structure de fibré en droites complexes sur $\mathcal{H} \otimes \mathbb{R}$
rendant holomorphes les sections locales de $\mathcal{H}$. Si $z_1, z_2$ sont deux
sections de $\mathcal{H}$ formant une base, la structure complexe équivaut à la
fonction holomorphe $\tau = z_2/z_1 : S \to \mathbb{C} \setminus \mathbb{R}$.

Une \emph{structure de Torelli} est un isomorphisme $t : \mathbb{Z}^2_S \simeq
\mathcal{H}$ ; la catégorie des courbes munies d'une telle structure est rigide,
et ses classes d'isomorphie correspondent aux applications holomorphes $\tau : S
\to \mathbb{C} \setminus \mathbb{R}$. La structure est \emph{stricte} si la base
$(z_1, z_2)$ est directe pour l'orientation définie par la structure complexe,
c'est-à-dire si $\tau$ est à valeurs dans $\mathfrak{D}^{+} = \{\Im \tau > 0\}$.
Ainsi $\mathfrak{D}^{+}$ est l'espace analytique de modules fin des courbes
munies d'une structure de Torelli stricte, avec la famille
\[
\mathfrak{X}_{\mathfrak{D}^{+}} = (\mathfrak{D}^{+} \times \mathbb{C}) /
\mathbb{Z}^2, \qquad T_{m,n}(\tau, z) = (\tau, z + m + n\tau),
\]
de fibres $\mathbb{C}/(\mathbb{Z} + \mathbb{Z}\tau) \simeq \mathbb{C}^{*}/q^{\mathbb{Z}}$,
$q = \exp(2i\pi\tau)$, $|q| < 1$.

$\operatorname{SL}(2, \mathbb{Z})$ opère à droite par $t \mapsto t \circ g$ :
$z'_1 = a z_1 + c z_2$, $z'_2 = b z_1 + d z_2$, d'où
\[
\tau \circ g = \frac{d\tau + b}{c\tau + a}, \qquad \tau \circ g^{-1} = \frac{a\tau -
b}{-c\tau + d} .
\]
Aucune de ces deux formules n'est l'action usuelle. Soit $\tau_\infty =
\operatorname{diag}(-1, 1) \in \operatorname{GL}(2, \mathbb{Z})$ ; la conjugaison
par $\tau_\infty$ change $\begin{pmatrix} A & B \\ C & D \end{pmatrix}$ en
$\begin{pmatrix} A & -B \\ -C & D \end{pmatrix}$, et $g \mapsto
\tau_\infty(g^{-1}) = \begin{pmatrix} d & b \\ c & a \end{pmatrix}$ --- « échange
de $a$ et de $d$ » --- est un anti-automorphisme, qui change l'action à droite
en une action à gauche :
\[
g \cdot \tau = \tau \circ \tau_\infty(g^{-1}) = \frac{a\tau + b}{c\tau + d} .
\]
C'est celle qu'on emploiera, « consacrée par l'usage ».\footnote{Les pages 146 et
147 sont rangées dans l'ordre inverse de celui du raisonnement : la page 147,
qui introduit $\tau_\infty$ (formules (11) à (13)), suit la page 145, et la page
146, qui porte la formule (14), lui fait suite. On lit dans l'ordre du
raisonnement. Le schéma de la page 146 porte sur sa flèche verticale un « 2 »
lu tel quel.} Au niveau de la famille, la multiplication par $(c\tau + d)^{-1}$
envoie $\mathbb{Z} + \mathbb{Z}\tau$ sur $\mathbb{Z} + \mathbb{Z}\, g\tau$, d'où
l'action
\[
T_g(\tau, z) = \Bigl( g\tau, \frac{z}{c\tau + d} \Bigr), \qquad T_{g,m,n}(\tau,
z) = \Bigl( g\tau, \frac{z}{c\tau + d} + m + n\, g\tau \Bigr)
\]
du produit semi-direct $\mathbb{Z}^2 \rtimes \operatorname{SL}(2, \mathbb{Z})$
sur $\mathfrak{D}^{+} \times \mathbb{C}$, dont on déduit celle de
$\operatorname{SL}(2, \mathbb{Z})$ sur $\mathfrak{X}_{\mathfrak{D}^{+}}$.

\begin{quote}
\textbf{Théorème} (page 148). La famille $\mathfrak{X}_{\mathfrak{D}^{+}} \to
\mathfrak{D}^{+}$, modulaire pour la rigidification de Torelli stricte, est un
revêtement universel de $M_{1,1}$, de groupe $\operatorname{SL}(2, \mathbb{Z})$
opérant par $g\tau = \frac{a\tau + b}{c\tau + d}$, c'est-à-dire par $\tau \circ
\tau_\infty(g^{-1})$.
\end{quote}
Les \emph{corollaires 1, 2, 3} (page 149) sont ceux de la première mouture :
$\mathcal{T}_{1,1} \simeq \operatorname{SL}(2, \mathbb{Z})$, et, par les
théorèmes sur les $\mathcal{T}_{g,\nu}$, le groupe des classes d'isotopie d'une
surface de genre $1$ pointée et celui des automorphismes extérieurs directs
d'un groupe de type $(1,1)$ s'identifient à $\operatorname{Aut}^{+}_{\mathbb{Z}}(H)$.

\emph{La rigidification différentielle} (pages 150 et 151). Une base $\varphi$
de $\underline{t}_E$ identifie $\underline{t}_E$ à $\mathbb{C}$, et $E$ à un
réseau de $\mathbb{C}$ ; avec en plus une structure de Torelli stricte, on a
$(z_1, z_2) \in \mathbb{C}^{*2}$ avec $z_2/z_1 \in \mathfrak{D}^{+}$, base directe
de $\mathbb{C}$. Sur
\[
P\mathfrak{D}^{+} = \{ (z_1, z_2) \mid z_2/z_1 \in \mathfrak{D}^{+} \}, \qquad
\mathfrak{X}_{P\mathfrak{D}^{+}} = (P\mathfrak{D}^{+} \times \mathbb{C})/\mathbb{Z}^2,
\]
\[
T_{m,n}(z_1, z_2, z) = (z_1, z_2, z + m z_1 + n z_2),
\]
opèrent $\operatorname{SL}(2, \mathbb{Z})$, par $T_g(z_1, z_2, z) = (dz_1 + cz_2,
bz_1 + az_2, z)$, et $\mathbb{C}^{*}$, par les dilatations $T_\lambda(z_1, z_2, z)
= \lambda^{-1}(z_1, z_2, z)$. Chacun opère librement, mais non leur produit :
$(-1, -1)$ agit trivialement sur $P\mathfrak{D}^{+}$, et par $z \mapsto -z$ sur les
fibres.\footnote{La fin de la page 151 est biffée de plusieurs traits ; on y lit
encore « d'ordre 2 ! », « section nulle » et, encadré, « $M_{1,1}$ ! ». La
page 152 est blanche. Le quotient $P\mathfrak{D}^{+}/\operatorname{SL}(2,
\mathbb{Z})$ est l'espace des réseaux de $\mathbb{C}$, que les invariants
$(g_2, g_3)$ identifient au complémentaire de la courbe $g_2^3 = 27g_3^2$ ;
comme $P\mathfrak{D}^{+} \simeq \mathbb{C}^{*} \times \mathfrak{D}^{+}$, son
groupe fondamental est une extension de $\operatorname{SL}(2, \mathbb{Z})$ par
$\mathbb{Z}$, le groupe de tresses $B_3$. C'est où menait le revêtement
universel de la page 140 ; la remarque est la nôtre.}

\pagerange{152}{156}
\section*{VI. Appendice : normalité des extensions quadratiques (pages 152 à 156)}

La page 153 s'ouvre sur un calcul de produits dans une extension $1 \to N \to G
\to \mathfrak{S}_A \to 1$, $N = \operatorname{Ker}(\mathbb{F}_2^A \to
\mathbb{F}_2)$, scindée par un relèvement de $\mathfrak{S}_A$ ; d'une autre
plume, sans lien apparent avec la suite.\footnote{Pour $|A| = 3$ c'est le
groupe $V_J(\mathbb{F}_2) \rtimes \mathfrak{S}_J \simeq \mathfrak{S}_4$ de la
partie I ; le rapprochement est le nôtre.} Suit, sous un trait, l'étude de la
normalité de $B = \mathbb{Z}[t][S]/(S^2 - Q(t))$ sur $A = \mathbb{Z}[t]$, en
codimension $1$.

\emph{Normalité de $B = A[S]/(S^2 - a)$} (page 154), pour $A$ local régulier de
dimension $1$ --- un anneau de valuation discrète, cas des anneaux locaux de
codimension $1$ auquel on l'applique ---, d'idéal maximal $\mathfrak{m}$ et de
caractéristique résiduelle $p$. La condition $a \notin \mathfrak{m}^2$ est
toujours nécessaire, et suffit si $p \neq 2$.\footnote{La page suppose $A$
« régulier » local, sans restriction de dimension. En dimension $\geq 2$
l'énoncé est faux : $k[[x, y]][S]/(S^2 - xy)$ est normal avec $xy \in
\mathfrak{m}^2$. En dimension $1$, si $a = \pi^2 b$, l'élément $S/\pi$ est
entier et n'est pas dans $B$.} Si $p = 2$ :
\begin{enumerate}
\item[a)] $a \in \mathfrak{m}$ : $B$ est normal si et seulement si $a \notin
\mathfrak{m}^2$ ;
\item[b)] $a$ inversible, d'image $\dot a$ dans $k = A/\mathfrak{m}$ :
$1^\circ$) si $\dot a \notin k^2$, $B$ est normal, de corps résiduel $k(\sqrt{\dot
a})$ ; $2^\circ$) si $\dot a = \dot\alpha^2$, écrivons $a = \alpha^2 - b$, $b \in
\mathfrak{m}$ ; avec $T = S - \alpha$, $B \simeq A[T]/(T^2 + 2\alpha T + b)$, et
comme $T^2$ et $2\alpha T$ sont dans le carré de l'idéal maximal ($2 \in
\mathfrak{m}$), $B$ est normal si et seulement si $b \notin \mathfrak{m}^2$.
\end{enumerate}
La classe de $b$ dans $\mathfrak{m}/\mathfrak{m}^2$ ne dépend pas du choix de
$\alpha$.\footnote{La note marginale se lit « ne dépend que du choix de
$\alpha$ », sous réserve. Changer $\alpha$ en $\alpha + m$, $m \in \mathfrak{m}$,
change $b$ de $2\alpha m + m^2 \in \mathfrak{m}^2$ : la classe est indépendante
du choix.} Les cas normaux sont donc : $a \in \mathfrak{m} \setminus
\mathfrak{m}^2$ ; $a = \alpha^2 - b$ avec $\alpha$ inversible et $b \in
\mathfrak{m} \setminus \mathfrak{m}^2$ ; $a$ inversible et $\dot a \notin k^2$.

\emph{Application aux revêtements quadratiques de $\mathbb{P}^1_{\mathbb{Z}}$}
(pages 154 et 155). La fibre générique est une extension quadratique de
$\mathbb{Q}(t)$, donnée par un élément de $\mathbb{Q}(t)^{*}/\mathbb{Q}(t)^{*2}$,
représenté de façon unique par $\pm q\, Q$ : $q$ entier positif sans facteur
carré, $Q \in \mathbb{Z}[t]$ primitif, sans facteur carré, de coefficient dominant
positif. Pour $p$ impair, l'extension n'est génériquement étale en $p$ que si
$p \nmid q$ ; donc $q \in \{1, 2\}$. En caractéristique $2$, $S^2 - P(T)$ n'est
jamais génériquement étale : pour obtenir quelque chose de génériquement étale
en $2$ après normalisation, il faut que $B(2)$ ne soit pas normal, d'où $q = 1$
et $\varepsilon Q$ carré modulo $4$ --- ce qui fixe le signe $\varepsilon$, $-1$
n'étant pas un carré modulo $4$. « Cette condition est aussi suffisante\ldots »,
affirme la page, sans le montrer.

\emph{Exemple} (« Legendre modulaires des courbes elliptiques ») : $Q(T) = T(T -
\lambda)$, $\lambda \in \mathbb{Z} \setminus \{0\}$, dans le cas qui intéresse
Grothendieck $\lambda = 2^\alpha 3^\beta$, « l'invariant modulaire
exceptionnel ».\footnote{On reconnaît $1728 = 2^6 \cdot 3^3$, valeur de $j$ en la
courbe à multiplication par $i$ : le revêtement double de la droite des $j$
ramifié en $j = 0$ et $j = 1728$. L'identification est la nôtre ; la page ne
nomme pas $j$. Elle écrit « $Q(T) \equiv T^2$ (2) si $\lambda \equiv 0$ (4) »,
où la condition utile est la congruence modulo $4$.} Si $4 \mid \lambda$, $Q
\equiv T^2 \pmod 4$, et l'on trouve un revêtement étale de $\mathbb{A}^1_{\mathbb{Z}}
\setminus \{0, \lambda\}$ ; la page passe à la coordonnée $T' = 1/T$ pour
regarder ce qui se passe à l'infini, et s'arrête sur la formule. La page 156
reprend quelques calculs ($T^2 + 2\alpha T + b$, $4T - 1$) dans les blancs d'une
convocation tamponnée des 9 et 16 juin 1981.

\end{document}
